%% paper.tex  ---  JPE submission draft
%% Diamond-Dybvig critique and replacement model
%% John S. Schuler, George Mason University
%%
%% Compile: knitr (if needed) -> pdflatex -> bibtex -> pdflatex x2

\documentclass[12pt]{article}\usepackage[]{graphicx}\usepackage[]{xcolor}
% maxwidth is the original width if it is less than linewidth
% otherwise use linewidth (to make sure the graphics do not exceed the margin)
\makeatletter
\def\maxwidth{ %
  \ifdim\Gin@nat@width>\linewidth
    \linewidth
  \else
    \Gin@nat@width
  \fi
}
\makeatother

\definecolor{fgcolor}{rgb}{0.345, 0.345, 0.345}
\newcommand{\hlnum}[1]{\textcolor[rgb]{0.686,0.059,0.569}{#1}}%
\newcommand{\hlsng}[1]{\textcolor[rgb]{0.192,0.494,0.8}{#1}}%
\newcommand{\hlcom}[1]{\textcolor[rgb]{0.678,0.584,0.686}{\textit{#1}}}%
\newcommand{\hlopt}[1]{\textcolor[rgb]{0,0,0}{#1}}%
\newcommand{\hldef}[1]{\textcolor[rgb]{0.345,0.345,0.345}{#1}}%
\newcommand{\hlkwa}[1]{\textcolor[rgb]{0.161,0.373,0.58}{\textbf{#1}}}%
\newcommand{\hlkwb}[1]{\textcolor[rgb]{0.69,0.353,0.396}{#1}}%
\newcommand{\hlkwc}[1]{\textcolor[rgb]{0.333,0.667,0.333}{#1}}%
\newcommand{\hlkwd}[1]{\textcolor[rgb]{0.737,0.353,0.396}{\textbf{#1}}}%
\let\hlipl\hlkwb

\usepackage{framed}
\makeatletter
\newenvironment{kframe}{%
 \def\at@end@of@kframe{}%
 \ifinner\ifhmode%
  \def\at@end@of@kframe{\end{minipage}}%
  \begin{minipage}{\columnwidth}%
 \fi\fi%
 \def\FrameCommand##1{\hskip\@totalleftmargin \hskip-\fboxsep
 \colorbox{shadecolor}{##1}\hskip-\fboxsep
     % There is no \\@totalrightmargin, so:
     \hskip-\linewidth \hskip-\@totalleftmargin \hskip\columnwidth}%
 \MakeFramed {\advance\hsize-\width
   \@totalleftmargin\z@ \linewidth\hsize
   \@setminipage}}%
 {\par\unskip\endMakeFramed%
 \at@end@of@kframe}
\makeatother

\definecolor{shadecolor}{rgb}{.97, .97, .97}
\definecolor{messagecolor}{rgb}{0, 0, 0}
\definecolor{warningcolor}{rgb}{1, 0, 1}
\definecolor{errorcolor}{rgb}{1, 0, 0}
\newenvironment{knitrout}{}{} % an empty environment to be redefined in TeX

\usepackage{alltt}

%% --- Packages ---
\usepackage[margin=1in]{geometry}
\usepackage{setspace}
\usepackage{amsmath,amssymb,amsthm}
\usepackage{natbib}
\usepackage{har2nat}
\usepackage{graphicx}
\usepackage{hyperref}
\usepackage{booktabs}



%% --- JPE section numbering ---
\renewcommand{\thesection}{\Roman{section}}
\renewcommand{\thesubsection}{\Roman{section}.\Alph{subsection}}

%% --- Theorem environments ---
\newtheorem{theorem}{Theorem}
\newtheorem{lemma}{Lemma}
\newtheorem{proposition}{Proposition}
\newtheorem{definition}{Definition}
\newtheorem{corollary}{Corollary}
\newtheorem{remark}{Remark}
\renewcommand{\proofname}{Proof}

%% --- Spacing and style ---
\doublespacing
\bibliographystyle{jpe}

%% --- Title block ---
\title{Deposits Are Not Options: Contract Artifacts, Equilibrium
       Selection, and Network Contagion in Bank Run Models}

\author{John S.\ Schuler\\
        \small Department of Computational and Data Sciences\\
        \small George Mason University\\
        \small \texttt{jschule4@gmu.edu}}

\date{}

%%============================================================
\IfFileExists{upquote.sty}{\usepackage{upquote}}{}
\begin{document}
%%============================================================

\maketitle

%%------------------------------------------------------------
\begin{abstract}
\citet{diamond1983bank} establishes that bank runs can arise as
coordination failures when a liquidity-insurance contract embeds an
above-par early-withdrawal premium. This paper argues that the
canonical DD panic mechanism is premium-amplified: run
incentives generated by the above-par premium channel are strictly
increasing in $\varepsilon = c_1 - 1$ in the formal model, while the
incremental shift in the marginal cutoff becomes small near zero.
Finite sequential service can nevertheless produce failure under any
positive premium when withdrawal pressure is sufficiently high.
Ordinary deposit contracts do
not contain this above-par redemption feature and therefore do not
reproduce the canonical optimal contract's marginal withdrawal
condition. This is a contract-structure mismatch, not a rejection of
DD's broader sequential-service logic: recovery priority and strategic
complementarity remain relevant in both environments. We formalize
the DD result using a stochastic process model with $K$ agents,
shifted CRRA utility, and random sequential activation. A
49{,}440-cell simulation separates a fixed-contract test from an ex
ante common-allocation and participation exercise. A bridge theorem
shows that as $\varepsilon
\to 0$ the depositor's problem ceases to be an above-par
option-exercise problem and can be represented as a
recovery-probability comparison. The theorem provides a formal
handoff from the canonical expected-utility problem to the decision
object used in
a network-contagion model in which depositors infer recovery risk
from local withdrawal observations under sequential service.
New simulations evaluate three decision rules---a relative-safety
comparison, the bridge threshold, and explicit expected utility---on
matched network realizations. Local information generates endogenous
withdrawals and bank failures under all three rules. The bridge threshold
and explicit-utility rules agree on 99.5\%
of common-scenario validation outcomes and
96.1\% of a prespecified
comparative-statics sweep. They diverge, as expected, in an adaptively
selected stress test reported in the appendix. Thus the network
mechanism is not an artifact of one decision rule, although the exact
rule matters near contested decision boundaries. Reserve ratios and
insurance coverage are measurable; the depositor observation
structure is potentially recoverable from suitable supervisory or
depositor-link data, although
identifying the economically relevant observation network remains an
empirical problem.
\end{abstract}

%%============================================================
\section{Introduction}
%%============================================================

\citet{diamond1983bank} established one of the most important results
in financial economics: bank runs can occur as a pure coordination
failure, even when the bank's underlying assets are sound. The model
has shaped regulatory thinking, deposit insurance policy, and the
theory of financial fragility for four decades.

We do not challenge the existence result. The question is which part
of its marginal incentive applies unchanged to ordinary deposits.
DD's optimal deposit contract provides liquidity insurance to
depositors by paying early withdrawers above the liquidation value
of the long asset---a premium above par. This premium is not incidental
to the canonical optimal contract. It enters directly into the
patient depositor's comparison between early withdrawal and
continuation. Ordinary demand deposits instead redeem at par, so they
do not reproduce that exact marginal payoff comparison. Sequential
service, recovery priority, and strategic complementarity nevertheless
remain common to both environments.

The mechanism DD identifies is amplified by this premium.
Conditional on the survival map in Theorem~\ref{thm:main}, the
withdrawal threshold rises and failure probability weakly increases
with $\varepsilon$. A first-order approximation
(Corollary~\ref{cor:approx}) shows that the withdrawal threshold
rises only modestly above its zero-premium limit at small premium
levels. The fixed-contract DD simulation
(Section~\ref{sec:ddsimulation}) confirms the aggregate upward premium
comparative static using independent holdout draws and locates the
liquidity boundary at high withdrawal pressure when the premium is
small. A separate common-allocation search measures participation
rather than treating nonformation as bank failure, while the
separate
finite-agent bridge simulation evaluates the limiting withdraw/stay
decision under sequential service. Real banking has no large premium: FDIC deposit insurance pays
par, not premium; early withdrawal penalties on time deposits move in
the opposite direction. The canonical above-par channel therefore
cannot be carried unchanged into an ordinary-deposit application.
What remains is the shared sequential-service problem: whether
recovery priority makes withdrawal preferable when waiting is risky.


This observation points to a methodological question that DD's
framework leaves open: what determines where between the two
equilibria the system actually goes. Answering this requires modeling
the marginal behavior of agents who are neither committed early
withdrawers nor committed holders, but are on the knife edge
responding to signals, premia, and, critically, the observed behavior of others.
Stochastic process techniques make that marginal behavior explicit.

The paper proceeds in three linked stages. First, we isolate the
above-par payment in DD's canonical optimal contract and establish its
effect on the marginal withdrawal threshold and run probability.
Second, we derive the limiting decision rule as that premium approaches
zero and evaluate the withdraw-versus-stay classification in a
finite-agent sequential-service environment with partial payments.
Third, we transfer the generalized recovery comparison into a network
setting in which depositors observe local withdrawals and form beliefs
about recovery risk. This architecture moves from the canonical model
to an ordinary par-value deposit environment without replacing
expected-utility behavior with an unrelated threshold assumption. The
contrast concerns the marginal payoff and information mechanisms, not
the relevance of sequential service itself.

The paper makes three contributions. First, it identifies a
contract-structure mismatch: the canonical DD marginal withdrawal
condition contains an above-par early-payment feature absent from
ordinary demand deposits. Second, it derives the limiting
recovery-probability rule that remains as this premium vanishes and
generalizes the rule to risky immediate withdrawal. Third, it embeds
that decision object in a local-information network with sequential
service and shows, conditional on the model and simulation designs,
that endogenous withdrawals and bank failures arise under comparative
recovery, a calibrated bridge threshold, and explicit expected
utility. The simulations additionally locate settings in which the
bridge approximation tracks explicit utility and selected
high-disagreement regions in which the exact rule is consequential.

The paper is related to several strands of literature.

\paragraph{The Diamond--Dybvig framework and its extensions.}
\citet{bryant1980model} first showed that a sequential-service
constraint generates bank fragility in a simple storage economy.
\citet{diamond1983bank} embedded this observation in a three-period
model with a productive long asset, an optimal deposit contract, and
two Nash equilibria. The panic equilibrium in DD exists precisely
because the optimal contract pays early withdrawers above the
liquidation value of the long asset; deposit insurance eliminates
it by making early withdrawal payoff-irrelevant for patient agents.
\citet{green2000diamond} showed that with history-dependent
contracts, the efficient allocation is achievable as a unique
equilibrium, challenging DD's claim that multiple equilibria are
an inherent feature of banking. \citet{peck2003equilibrium}
re-established multiplicity in a model in which the bank cannot
commit to a payment rule ex ante, showing that the strategic
interaction in the withdrawal queue survives even under deposit
insurance when commitment is imperfect. \citet{ennis2010fundamental}
trace this fragility more broadly to the bank's inability to commit
to a payment rule; the same commitment problem reappears in
Section~\ref{sec:policy}, where it limits what a discretionary policy
response can credibly promise. \citet{jacklin1987demand} showed that
a demand-deposit contract combined with restricted trading provides
better risk sharing than a freely traded equity claim, establishing
why the demand-deposit form---rather than a generic risk-sharing
arrangement---is the relevant object to which the above-par premium
mechanism should be compared. Our contribution to this
strand is not to contest the existence results but to identify the
contractual feature---the above-par premium---that changes the
canonical optimal contract's marginal withdrawal condition and is
absent from ordinary par-value demand deposits.

\paragraph{Equilibrium selection and global games.}
The indeterminacy of equilibrium selection in DD---which of two
equilibria obtains is determined by an exogenous sunspot---motivated
a literature on equilibrium selection through information
heterogeneity. \citet{morris1998unique} showed that small amounts
of heterogeneous private information eliminate multiplicity in
coordination games, yielding a unique threshold equilibrium.
\citet{goldstein2005demand} applied this global-games methodology
directly to bank runs, deriving a unique equilibrium in which
depositors withdraw when their private signal about asset quality
falls below a threshold $p^*$. \citet{rochet2004coordination}
extended the framework to interbank coordination failures. The
global-games threshold is formally related to the threshold
$p^*(\rho, R)$ of our bridge theorem, but the economic content
and information structure differ. In Goldstein and Pauzner,
$p^*$ is a cutoff on a private signal about the bank's fundamental
soundness; in our model it is the limit of rational CRRA expected
utility as the insurance premium vanishes, and agents hold no
private fundamental signal---they form beliefs from locally
observed withdrawal behavior. The two frameworks are better understood as complementary than
competing, because they model different information environments.
Global games apply when depositors hold private signals about bank
fundamentals---asset quality, reserve adequacy, managerial
competence---and coordinate on a unique equilibrium through the
noise in those signals. Our model applies when depositors lack
such private fundamental information and instead rely entirely on
observation from their local network. SVB depositors held standard
uninsured commercial deposits with no private window into the bank's
bond portfolio. \citet{cookson2023social} report that measures of
social-media activity in SVB-connected venture-capital and startup
networks precede and predict deposit outflows in Granger-causality
tests. This temporal predictive relationship is consistent with
depositors responding to information circulating through their
networks, but it does not by itself identify a structural causal
channel. \citet{iyer2012understanding} provide complementary,
non-crisis evidence for the same premise using depositor-level Indian
bank data: depositors who observe a co-depositor's early withdrawal
are more likely to withdraw themselves, and this network effect is
distinct from and additive to fundamentals-based inference. The two
information structures are empirically
distinguishable and call for different models.

The distinct information structures correspond to distinct
predictions. In standard global-games applications, run risk is
governed by fundamentals and private-signal precision. In our model,
conditional run risk also varies with the depositor observation
network. The baseline assumes that network position is independent
of deposit size; this restriction is testable with depositor-level
supervisory data. Information timing, sequenced redemption, and
temporary suspension also act directly on the observability margin in
our model. Such interventions can be represented in richer
global-games environments---\citet{angeletoswerning2006crises} show
that endogenous public information from prices can itself generate
multiplicity and excess volatility in a global-games setting, so the
literature does have a policy margin operating through signal
precision and publicity. That margin is nonetheless distinct from
ours: it operates on the precision or publicity of an exogenous
private signal, whereas our margin operates on whether an endogenous
individual action---a withdrawal---is observable at all, and to whom.
The two are related but not the same lever, and neither is absent
from the literature.

\paragraph{Information-based runs and rational inference.}
A parallel strand treats bank runs as consequences of rational
inference from partial signals. \citet{jacklin1988distinguishing}
distinguished information-based runs, in which withdrawals reflect
genuine adverse information about asset quality, from panic-based
runs driven purely by coordination failure. \citet{chari1988banking}
constructed an equilibrium in which depositors observe the
withdrawal queue length and update beliefs about asset quality,
producing runs as rational cascades rather than sunspot phenomena.
\citet{hemanela2016information} model a related but distinct
inference channel: depositors decide whether to pay a cost to
verify a rumor about bank solvency, so information acquisition
itself is endogenous and strategic. Our mechanism does not involve
costly verification or rumors about fundamentals; agents instead
observe realized withdrawal decisions directly from network
neighbors at no cost, and the object they are uncertain about is the
unobserved withdrawal count, not asset quality. Our model is closest to this literature and extends the inference logic to a network setting:
agents observe withdrawal rates among network neighbors and
update beliefs about the global withdrawal distribution. The
mechanism---extrapolation from a partial signal
producing cascades---is in this tradition. What is novel is the
information structure (local network sampling rather than queue
position) and its implication that cascade susceptibility depends
on neighborhood degree through a local-sample-quality channel and on
shortcut density through information transmission.

\paragraph{Financial networks and cascade models.}
\citet{allen1998optimal} characterized efficient responses to
aggregate liquidity shocks in a banking system and showed that
market incompleteness can generate inefficient crises.
\citet{allen2000financial} introduced balance-sheet network
contagion, showing that interbank credit linkages propagate
localized shocks across the system and that incomplete networks
are more contagion-prone than complete ones under moderate shocks.
\citet{acemoglu2015systemic} established a phase transition:
denser financial networks are more stable under small shocks but
amplify large ones. Our observation network is informational
rather than financial---agents are linked by mutual observation
of withdrawal behavior, not by credit claims---so the contagion
mechanism is belief propagation rather than loss propagation.
This places our model closer to \citet{watts2002simple}'s
threshold cascade framework, in which nodes adopt a behavior
when the fraction of their neighbors doing so exceeds a threshold.
Our withdrawal rule---withdraw when the locally estimated loss
probability conditional on staying exceeds that conditional on withdrawal---
is similar to the threshold cascade
applied to a belief conditional on prior local withdrawals, on a Newman--Watts--Strogatz
network chosen to capture the small-world topology of social and
professional depositor communities \citep{newman1999renormalization}.

\paragraph{Social media, digital banking, and the SVB run.}
The behavior of wholesale and shadow-banking instruments during
recent crises motivates the scope claims of
this paper. \citet{gorton2012securitized} documented the 2007--2008
run on repo: overnight repurchase agreements are redeemable at
par on demand and experienced classic runs when counterparties
refused to roll over, consistent with DD's mechanism applying
to instruments that approximate above-par early redemption.
Money market fund runs in 2008 and 2020 share this structure.
The SVB run of March 2023 contrasts sharply: SVB's depositors
held standard uninsured commercial deposits with no above-par
redemption feature. \citet{cookson2023social} report that measures
of Twitter activity among SVB-connected venture-capital firms and
startups precede and predict deposit outflows in Granger-causality
tests. The result is consistent with propagation through an
informational observation network, but does not establish that
network propagation caused the outflows.
The speed of the SVB run---withdrawals occurring at very high
frequency---is more naturally accommodated by an information-cascade
account than by a literal above-par option-exercise interpretation.
This is a comparison of mechanism fit, not identification from timing
alone. SVB therefore motivates the contract-structure distinction and
the study of informational propagation on depositor networks.

\paragraph{Agent-based and stochastic-process approaches.}
The limitations of closed-form equilibrium analysis in capturing
run dynamics---path dependence, activation order, and marginal
agent behavior between equilibria---have motivated a growing
literature using agent-based and stochastic-process methods
\citep{bookstaber2017end}. These approaches make the trajectory
between equilibria explicit. Our simulation implements the network-contagion model as a
sequential-activation agent-based model in this tradition, but
with an explicit microfoundation motivated by the bridge theorem.


The paper proceeds as follows. Section~\ref{sec:dd} presents the
Diamond--Dybvig model and derives the upward slope result.
Section~\ref{sec:process} develops the stochastic process model and
proves the main formal theorems. Section~\ref{sec:bridge} presents
the bridge theorem. Section~\ref{sec:replacement} introduces the
network-contagion model. Section~\ref{sec:results} presents network
comparative statics from the agent-based model. Section~\ref{sec:policy}
contrasts the two models and develops policy implications.
Section~\ref{sec:conclusion} concludes.

%%============================================================
\section{The Diamond--Dybvig Model and the Option Premium}
\label{sec:dd}
%%============================================================

\subsection{Setup}
 


We follow the standard DD setup. There are three periods
$t \in \{0, 1, 2\}$ and $K$ agents, each endowed with one unit at
$t=0$. Agents are ex ante identical but face private liquidity shocks
at $t=1$: with probability $\pi$ an agent is type~1 (impatient) and
values only $t=1$ consumption; with probability $1-\pi$ an agent is
type~2 (patient) and values only $t=2$ consumption. Types are private
information.

A long asset pays gross return $R > 1$ at $t=2$ but yields only $1$
if liquidated at $t=1$. A short asset (storage) yields $1$ in any
period.

The bank accepts deposits at $t=0$ and offers a deposit contract
$(c_1, c_2)$ where $c_1$ is the $t=1$ withdrawal payment and $c_2$
is the $t=2$ withdrawal payment. Subject to a sequential service
constraint, the bank pays $c_1$ to withdrawers at $t=1$ until funds
are exhausted.

\subsection{The Optimal Contract and the Resource Constraint}

The social planner maximizes ex ante expected utility:
%
\begin{equation}
  \max_{c_1,\,c_2} \; \pi\, u(c_1) + (1-\pi)\, u(c_2)
  \label{eq:planner}
\end{equation}
%
subject to the resource constraint. If $n$ denotes the fraction of
agents withdrawing at $t=1$, the resource constraint requires:
%
\begin{equation}
  n c_1 + (1-n) c_2 \;\leq\; n c_1 + \bigl(1 - n c_1\bigr) R
  \label{eq:resource_full}
\end{equation}
%
which simplifies to:
%
\begin{equation}
  (1-n)\, c_2 \;\leq\; \bigl(1 - n c_1\bigr) R
  \label{eq:resource}
\end{equation}

The optimal contract sets $c_1 > 1$. This is the liquidity insurance
premium---early withdrawers receive strictly more than the storage
value of liquidated assets. The premium $\varepsilon \equiv c_1 - 1 > 0$
is the option value of early withdrawal.

The resource constraint ties $c_1$ and $c_2$ together: for fixed $n$
and $R$, raising $c_1$ necessarily depresses $c_2$.

\subsection{The Two Equilibria}

DD establishes two Nash equilibria.

\paragraph{Fundamental equilibrium.} Only type~1 agents withdraw at
$t=1$. The bank is solvent. Patient agents receive $c_2 > c_1$.

\paragraph{Panic equilibrium.} All agents withdraw at $t=1$. The bank
exhausts funds before serving all depositors. Patient agents who
withdraw late receive nothing.

The panic equilibrium exists because $c_1 > 1$: when a patient agent
believes sufficiently many others will withdraw, early withdrawal
dominates waiting. The above-par premium is the mechanism specific to the DD optimal
liquidity-insurance contract. At $c_1 = 1$, the contract loses the
above-par feature that makes early withdrawal option-like for patient
agents: the bank offers no improvement over storage and agents do not
bank under the DD contract. This rules out the DD above-par premium
channel as the operative mechanism---it does not rule out runs in
broader sequential-service environments where early withdrawal pays
par but waiting is risky. The distinction is between an above-par
option-like payoff and a par-now-versus-risky-later first-mover
advantage. The former is the DD optimal-contract mechanism; the
latter is a broader run mechanism that can operate in ordinary
deposits and motivates the network-contagion model below.

\subsection{The Upward Slope: Withdrawal Incentive Is Increasing in $c_1$}

DD establishes that the panic equilibrium exists for $c_1 >
1$ but says nothing about the probability of reaching it.
A necessary first step is establishing that the individual incentive
to withdraw is itself increasing in $c_1$: as the premium rises,
each agent requires a higher assessed survival probability to prefer
waiting over withdrawing.

\begin{proposition}
\label{prop:upslope}
Under the DD resource constraint~\eqref{eq:resource} and shifted CRRA
utility, the withdrawal threshold
$\Phi_{\rm shift}(c_1,\,c_2(c_1),\,\rho)$ is strictly increasing in
$c_1$ for all $\rho > 0$.
\end{proposition}

The mechanism operates through two simultaneous channels. Raising
$c_1$ increases the value of early withdrawal directly. It also
depresses $c_2$ through the resource constraint, reducing the value
of waiting. Both channels increase $\Phi_{\rm shift}$, enlarging the
set of survival beliefs under which withdrawal is individually
rational. The proof is in Appendix~\ref{app:proofA} (Lemma~\ref{lem:phi}).

This is a purely formal result: it follows from the utility
specification and the resource constraint alone, with no reference to
equilibrium dynamics or the stochastic withdrawal process.
The translation from a higher withdrawal threshold to a higher
probability of the panic equilibrium requires the stochastic
process model of Section~\ref{sec:process}, where a third channel
also operates: larger $c_1$ increases the jump size of each
withdrawal event, accelerating pool depletion. The full monotonicity
result for run probability is Theorem~\ref{thm:main}.

\subsection{The Scope Condition}

The premium $\varepsilon=c_1-1>0$ defines a scope condition for the
canonical DD optimal contract's exact marginal withdrawal condition.
It is absent from ordinary demand deposits, which redeem at par.
This does not imply that $c_1>1$ is necessary for every run
equilibrium. Sequential service can create recovery-priority and
strategic-complementarity incentives even when $c_1=1$. The narrower
claim is that replacing $c_1>1$ with par redemption changes the
patient depositor's marginal payoff comparison and removes the
above-par component of the canonical mechanism.

\paragraph{Quantitative smallness.} The stronger claim concerns the
marginal cutoff: even where $\varepsilon > 0$, a small premium
contributes only a modest increment to that cutoff
at ordinary-deposit premium levels. FDIC deposit
insurance pays par, not premium; early withdrawal penalties on time
deposits are below par; and the instruments most
closely resembling DD's contract---money market funds, repo agreements,
certain wholesale funding structures---carry premia far below the
level at which the option channel dominates. Corollary~\ref{cor:approx}
below establishes that the withdrawal threshold $\Phi_{\rm shift}$
exceeds its zero-premium limit by only a small amount at the premiums
found in ordinary deposit contracts. The result isolates the
incremental premium channel; it does not eliminate the remaining
recovery-priority incentive at the zero-premium limit. The instruments
that experienced runs in 2008 and 2020 are
discussed further in Section~\ref{sec:policy}.


The regulatory record documents this quantitative claim. All time
deposits in the United States must impose a minimum early withdrawal
penalty as a condition of their regulatory classification as time
deposits \citep{reg_d}; the Truth in Savings Act requires depository
institutions to disclose penalty terms before account opening
\citep{reg_dd}. Standard retail certificates of deposit carry
penalties of 90 to 180 days' simple interest for maturities of one
to two years, translating to 0.9 to 1.9 percent of principal at
current certificate rates---below par throughout. Demand deposits and
transaction accounts carry no early withdrawal premium or penalty by
definition ($\varepsilon = 0$ exactly). FDIC deposit insurance
reimburses insured balances at par by statute \citep{fdia_1821f},
so $\varepsilon = 0$ for all insured balances regardless of
instrument type. No standard
US retail deposit product pays an above-par early withdrawal premium.
The wholesale instruments most closely approximating DD's
contract---money market funds and overnight repo---experienced runs
in 2008 and 2020, but those runs took the form of rollover refusal
and gate imposition rather than above-par redemption
\citep{gorton2012securitized}, a distinction developed further in
Section~\ref{sec:policy}. Even treating $\varepsilon = 0.02$ as a
generous upper bound for any instrument with a positive premium,
Corollary~\ref{cor:approx} implies the withdrawal threshold rises
only modestly above its zero-premium limit. This is not a uniform
bound on failure probability: under sequential service, sufficiently
widespread withdrawals can exhaust a fixed pool for any
$\varepsilon>0$.

Uninsured deposits introduce a different object. Their contractual
claim may still redeem at par, so $c_1=1$ and $\varepsilon=0$, while
the realized payment can be partial or zero if reserves are depleted
before the claim is served. This is recovery risk, not a negative
contractual premium. It enters the depositor's problem through the
probability and distribution of payment under each action. In the
binary limiting comparison developed below, those objects are $p_W$,
the probability of full recovery after withdrawing, and $p_S$, the
probability of full recovery after staying. The March~2023 SVB run
involved uninsured par claims with uncertain recovery, making
recovery priority salient without placing $\varepsilon$ below zero.
This distinction motivates the generalized bridge condition rather
than an extension of the DD premium comparative static into
$\varepsilon<0$.

The panic equilibrium in DD therefore has limited relevance
for ordinary retail and commercial deposits: the mechanism is a
premium-amplified mechanism, whereas ordinary deposits operate at par
or with penalties. DD's existence proof stands; what is limited is
the direct transfer of its above-par option-exercise margin to the
most common form of banking relationship. Recovery priority under
sequential service remains operative at par.

%%============================================================
\section{Diamond--Dybvig Stochastic Process Model}
\label{sec:process}
%%============================================================

\subsection{Setup}

Each agent deposits their full endowment at $t=0$ under the optimal
contract of Section~\ref{sec:dd}, which specifies the payment pair
$(c_1, c_2)$. This section takes the contract as given and analyzes
the $t=1$ withdrawal timing decision; the deposit level is not a
choice variable here.

There are $K$ agents. The deposit pool $D(t) \in \mathbb{R}_{\geq 0}$
evolves as a pure jump process: each period, agents are processed in a
uniformly random sequential order; when an agent withdraws, $D(t)$
decrements by $c_1$. Exogenous (liquidity-shock) withdrawals are
determined by a single Binomial draw at the start of each run;
endogenous withdrawal decisions follow in the resulting random queue.

For consumption outcomes that may equal zero, agents have normalized
shifted CRRA utility:
%
\begin{equation}
  u_a(c) =
  \frac{(a+c)^{1-\rho}-a^{1-\rho}}{1-\rho},
  \quad \rho > 0,\ \rho \neq 1;
  \qquad
  u_a(c) = \log(a+c)-\log a,\quad \rho=1,
  \quad a>0.
  \label{eq:crra}
\end{equation}
%
The baseline calibration sets $a=1$. The positive shift avoids the
singularity at zero consumption, but it is not interchangeable with
unshifted CRRA: because it changes utility differences, it can change
the withdrawal cutoff. The sensitivity exercise therefore varies
$a\in\{0.1,1,10,100\}$ and reports the shift as a distinct
calibration dimension.

Agents are fully rational expected utility maximizers. The Stochastic DD Simulation
simulation computes the expectation in the withdrawal
condition~\eqref{eq:wdcond} numerically via Monte Carlo; this is a
computational implementation of rational calculation, not a
bounded-rationality approximation. The qualitative
results---monotonicity of run probability in $\varepsilon$
(Theorem~\ref{thm:main}) and convergence of the withdrawal condition
to a recovery-probability threshold as $\varepsilon \to 0$
(Theorem~\ref{thm:bridge})---are robust to the
specific endowment normalization used in the simulation.

\subsection{The Withdrawal Decision}

An agent activated when the deposit pool stands at $D$ solves:

\begin{quote}
Withdraw now: receive $c_1$ with certainty. \\[2pt]
Wait: receive $c_2$ with probability
$P(\text{survival to } t=2 \mid D)$, and $0$ otherwise.
\end{quote}
%
The agent withdraws if and only if:
%
\begin{equation}
  u(c_1) \;\geq\;
  P(\text{survival} \mid D)\cdot u(c_2) \;+\;
  \bigl[1 - P(\text{survival} \mid D)\bigr]\cdot u(0)
  \label{eq:wdcond}
\end{equation}
%
Rearranging~\eqref{eq:wdcond} and using $u(c_2)-u(0)>0$ for all
$\rho>0$ and $c_2>0$:
%
\begin{equation}
  P(\text{survival} \mid D) \;\leq\;
  \frac{u(c_1)-u(0)}{u(c_2)-u(0)}
  \;\equiv\; \Phi_{\rm shift}(c_1,\,c_2,\,\rho)
  \label{eq:phi}
\end{equation}
%
With shifted CRRA this equals $[(1+c_1)^{1-\rho}-1]/[(1+c_2)^{1-\rho}-1]$
for $\rho\neq 1$ and $\log(1+c_1)/\log(1+c_2)$ for $\rho=1$
(see equation~\eqref{eq:phishiftA} and Appendix~\ref{app:proofB}).
These displayed expressions set $a=1$. For general $a$, replace
$1+c$ by $a+c$ and the subtracted $1$ by $a^{1-\rho}$; for log
utility use differences from $\log a$.
The agent withdraws when the bank's survival probability falls below
$\Phi_{\rm shift}$. Two comparative statics follow directly
from~\eqref{eq:phi}: $\Phi_{\rm shift}$ is strictly increasing in
$c_1$ for all $\rho>0$ (Lemma~\ref{lem:phi}), and under the resource
constraint~\eqref{eq:resource} raising $c_1$ also decreases $c_2$,
which further increases $\Phi_{\rm shift}$. Both effects push toward
withdrawal.

\subsection{The Rational Expectations Equilibrium}

In rational expectations equilibrium, the withdrawal threshold $D^*$
must be consistent with the aggregate process it generates.

\begin{definition}
A rational expectations equilibrium is a threshold
$D^* \in [0,\, K c_1]$ such that an agent activated when $D = D^*$
is exactly indifferent between withdrawing and waiting, given that
all other agents use the same threshold $D^*$.
\end{definition}

The fixed point is characterized by:
%
\begin{equation}
  P\!\left(\tau_{D^*} \geq T \;\middle|\; D(0) = D^*\right)
  \;=\; \Phi(c_1,\,c_2,\,\rho)
  \label{eq:fp}
\end{equation}
%
where $\tau_{D^*}$ is the first time $D(t)$ hits $0$ when all agents
with $D \leq D^*$ withdraw, and $T$ is the maturity date. Thus
$\{\tau_{D^*}\geq T\}$ is the survival event. Let
$S(D;d,c_1)$ denote survival from pool $D$ when all agents use
threshold $d$, and define the equilibrium composite
$\bar S(d;c_1)=S(d;d,c_1)$. We impose two regularity conditions:
(R1) $\bar S$ is continuously differentiable, rises from 0 to 1 as
$d$ increases from 0 to $Kc_1$, and satisfies
$\partial\bar S/\partial d>0$ on the interior; and
(R2) $\Phi(c_1,c_2,\rho) < 1$, which holds
whenever $c_1 < c_2$---guaranteed by the resource constraint when
$R > 1$ and $\pi < 1$. Under (R1)--(R2), a solution
to~\eqref{eq:fp} exists and is unique for that survival map.

\subsection{Main Theorem}

\begin{theorem}[Conditional comparative static]
\label{thm:main}
Given an equilibrium survival composite $\bar S(d;c_1)$ satisfying
(R1)--(R2), and under the DD resource constraint~\eqref{eq:resource},
the rational expectations equilibrium
threshold $D^*$ is strictly increasing in $c_1$. The probability
of failure before maturity, $P(\tau < T)$, is weakly increasing in
$c_1$. It is strictly increasing when raising $c_1$ moves a
positive-probability set of coupled paths across the failure
boundary.
\end{theorem}
\begin{proof}[Proof sketch]
Raising $c_1$ operates through three reinforcing channels.
\begin{enumerate}
  \item Direct option effect. $\Phi$ is strictly increasing
    in $c_1$ (equation~\eqref{eq:phi}).
  \item Continuation value effect. Under
    constraint~\eqref{eq:resource}, $c_2$ is strictly decreasing in
    $c_1$, which further increases $\Phi$.
  \item Jump-size effect. Each withdrawal event reduces $D(t)$
    by $c_1$; larger jumps accelerate pool depletion for any fixed
    $D^*$.
\end{enumerate}
Channels~1 and~2 increase $\Phi$. The larger withdrawal payment also
weakly lowers survival at any fixed pool and withdrawal region.
Because the survival side of~\eqref{eq:fp} is strictly increasing in
the equilibrium pool under (R1), the implicit function theorem then
implies a higher $D^*$. A higher $D^*$ enlarges the withdrawal
region, increasing $P(\tau<T)$. In addition, Channel~3 increases
$P(\tau<T)$ for any fixed $D^*$ by the same coupling argument.
Full details are in Appendix~\ref{app:proofA}.
\end{proof}

\begin{remark}[Formal model and computational operationalization]
\label{rem:partial}
Theorem~\ref{thm:main} is conditional on (R1)--(R2), not a derivation
of the equilibrium survival composite from primitive withdrawal hazards. Continuity,
strict monotonicity, and therefore uniqueness are properties imposed
on $\bar S(d;c_1)$. The formal result establishes the comparative static
for any map in that class; it does not claim that every stochastic
withdrawal law belongs to the class or pin down $P(\tau<T)$ for a
given $(c_1,K,p_0)$. The fixed-contract DD simulation supplies
a finite-agent distributional specification; the separate allocation
panel evaluates a selected common allocation on independent holdout
draws. The distinct bridge
simulation supplies paired withdraw-now and stay counterfactuals
under sequential service. Together they operationalize the premium
comparative static and the limiting decision object without treating
the network model as a numerical proof of the theorem.
\end{remark}

The next section evaluates this DD comparative static before the
paper turns to the bridge theorem. Keeping the formal result and its
simulation adjacent separates the canonical-contract evidence from
the later bridge and network exercises.

%%============================================================
\section{Diamond--Dybvig Simulation}
\label{sec:ddsimulation}
%%============================================================

The simulation separates two questions that should not be conflated.
The theorem-facing experiment fixes the common deposit at the full
1{,}000-unit endowment and evaluates 41{,}600 parameter cells on
10{,}000 independent realizations each. The grid spans eight premia
from zero to 0.50, four productivity values, 100 prior withdrawal
probabilities, and thirteen nonredundant risk-aversion/utility-shift
specifications. The separate allocation experiment searches
$d\in\{0,10,\ldots,1000\}$ using 1{,}000 realizations per candidate
and evaluates the selected allocation on 10{,}000 holdout
realizations in each of 6{,}400 cells. This is a common ex ante
allocation choice, not a decentralized choice by heterogeneous
depositors. A selected deposit of zero means that no bank is formed;
it is nonparticipation rather than insolvency. Common random numbers
match premium and utility-shift comparisons, while allocation search
and holdout streams remain separate.

Failure is defined as a contractual early-payment shortfall. Exact
exhaustion after every claim has been honored is settlement, not
default. This distinction is decisive at the zero-premium endpoint:
all 5{,}200 zero-premium fixed-contract cells settle without default,
including $P_w=1$. At $P_w=1$, every positive-premium cell defaults
because the fixed pool cannot pay an above-par amount to every
depositor.

Figure~\ref{fig:ddcore_new} reports the baseline $a=1$
fixed-contract results. Averaged over the grid, failure rises from
5.8\%
at a premium of 0.005 to
67.7\%
at a premium of 0.50. The median prior-withdrawal probability at
which failure first exceeds one half moves from approximately 0.95
to 0.35 over the same range. Higher productivity is stabilizing
throughout the matched sweep. Premium monotonicity is strong in
aggregate but not literally pointwise: nine of 5{,}200 matched
premium paths have a simulated decline exceeding one percentage
point. These reversals occur near the high-withdrawal boundary and
mainly under the extreme shift $a=100$. The evidence therefore
supports an aggregate and boundary comparative static, not the claim
that every finite Monte Carlo cell is monotone.

\begin{figure}[ht]
\centering
\begin{knitrout}
\definecolor{shadecolor}{rgb}{0.969, 0.969, 0.969}\color{fgcolor}
\includegraphics[width=0.88\textwidth]{figures/knitr-dd-core-figure-1} 
\end{knitrout}
\caption{Fixed-contract Diamond--Dybvig simulation at utility shift
$a=1$. Each curve averages over four productivity values. Deposits
are fixed at the full endowment and each parameter cell is evaluated
on 10{,}000 realizations. Zero premium is exactly serviceable; higher
premia move the failure boundary toward lower prior withdrawal
probabilities.}
\label{fig:ddcore_new}
\end{figure}

The allocation search produces a complementary participation result.
Participation falls from
97.9\%
at zero premium to
33.0\%
at a premium of 0.50, while the mean selected deposit falls from
960.4 to
317.1.
The zero-premium nonparticipation cells are concentrated at
$P_w=1$, where outside storage and deposit followed by par withdrawal
are payoff-equivalent; they are tie outcomes, not evidence that the
banking contract is strictly rejected.

\begin{figure}[ht]
\centering
\begin{knitrout}
\definecolor{shadecolor}{rgb}{0.969, 0.969, 0.969}\color{fgcolor}
\includegraphics[width=0.88\textwidth]{figures/knitr-dd-allocation-figure-1} 
\end{knitrout}
\caption{Participation and selected deposits in the common-allocation
search. Each point averages over four productivity values and fifty
prior withdrawal probabilities. Both bank formation and the selected
deposit generally decline as the contractual premium rises.}
\label{fig:ddallocation_new}
\end{figure}

The option channel does not generally sustain maximum deposits in
fragile cells. Among the 142 participating
cells with failure rates of at least 10 percent, selected deposits
range from 10 to
1000. Under nonlinear utility,
none of these high-failure cells selects the full deposit. Deposit
choice is also not pointwise monotone in the premium because the
allocation jointly trades outside storage, early-payment option
value, late returns, and run exposure.

Finally, the utility shift is mathematically part of preferences but
numerically modest for the participation conclusion. In the matched
four-shift comparison, participation agrees in every cell; exact
deposit agreement is at least
85.0\%
in every risk-aversion/premium group, and the largest within-cell
deposit range is 100 units out
of 1{,}000. Appendix~\ref{app:ddsweep} reports the expanded grid and
shift sensitivity. Thus $a$ is not an innocuous algebraic
normalization, but the participation result is robust over the
calibrated range.

%%============================================================
\section{The Bridge Theorem}
\label{sec:bridge}
%%============================================================

The bridge theorem is best understood as a translation between
models. Its purpose is not to treat a continuity calculation as the
paper's main technical achievement. It identifies, within the
canonical expected-utility problem familiar to readers, the decision
object that survives when the above-par premium vanishes. That object
can then be carried into a sequential-service environment with risky
withdrawal and local information. The explicit threshold below and
its first-order expansion provide the formal handoff; the subsequent
simulations evaluate its finite-agent and network consequences.

\subsection{Statement}

The CRRA withdrawal condition~\eqref{eq:phi} requires agents to
compute $\Phi(c_1, c_2, \rho)$---a ratio that depends on the premium,
the continuation value, and the risk aversion parameter. We show that
its zero-premium limit can be represented as a comparison against a
fixed threshold.

Under the resource constraint~\eqref{eq:resource} at the fundamental
equilibrium ($n = \pi$):
%
\begin{equation}
  c_2(\varepsilon) \;=\;
  \frac{\bigl(1 - \pi(1+\varepsilon)\bigr)R}{1-\pi}
  \;=\; R - \frac{\pi\varepsilon R}{1-\pi}
  \label{eq:c2eps}
\end{equation}
%
so $c_2(0) = R$ and $dc_2/d\varepsilon\big|_0 = -\pi R/(1-\pi) < 0$.

\begin{theorem}[Bridge Theorem]
\label{thm:bridge}
Let $\varepsilon = c_1 - 1 \geq 0$. Under the shifted CRRA utility
specification and resource constraint~\eqref{eq:resource}, as
$\varepsilon \to 0$ the withdrawal condition converges to:
%
\begin{equation}
  P(\textup{survival} \mid D) \;\leq\; p^*(\rho,\, R)
  \label{eq:bridge}
\end{equation}
%
where the threshold is:
%
\begin{equation}
  p^*(\rho,\, R) \;=\;
  \frac{2^{1-\rho} - 1}{(1+R)^{1-\rho} - 1},
  \quad \rho \neq 1;
  \qquad
  p^*(1,\, R) \;=\; \frac{\log 2}{\log(1+R)}
  \label{eq:pstar}
\end{equation}
The function $p^*$ is continuous at $\rho = 1$. For $\rho \neq 1$,
$p^*$ depends on both $\rho$ and $R$; at $\rho = 1$ (log utility),
L'H\^opital's rule eliminates $\rho$ from the expression and $p^*$
is pinned by $R$ alone. This is not a claim that preferences are
irrelevant---log utility is itself a preference specification---but
that, conditional on log utility, no residual risk-aversion parameter enters
the threshold. The Stochastic DD Simulation fixes $\rho = 1$ for this reason:
under log utility the decision threshold is fully determined by $R$,
so the simulation need not sweep over $\rho$ as an additional free
parameter.
\end{theorem}

\begin{proof}[Proof sketch]
With $c_1 = 1 + \varepsilon$ and $c_2(\varepsilon)$ given
by~\eqref{eq:c2eps}, the shifted CRRA withdrawal condition
(see Appendix~\ref{app:proofB}) rearranges to
$P \leq \Phi_{\rm shift}(\varepsilon)$ where:
\[
  \Phi_{\rm shift}(\varepsilon) \;=\;
  \frac{(1+c_1)^{1-\rho} - 1}{(1+c_2)^{1-\rho} - 1}
  \;=\;
  \frac{(2+\varepsilon)^{1-\rho} - 1}
       {\bigl(1 + c_2(\varepsilon)\bigr)^{1-\rho} - 1}
\]
Since $c_1 \to 1$ and $c_2 \to R$ as $\varepsilon \to 0$, continuity
gives $\Phi_{\rm shift}(\varepsilon) \to p^*(\rho, R)$ as stated.
For $\rho = 1$, the corresponding condition is
$P \leq \log(1+c_1)/\log(1+c_2) \to \log 2/\log(1+R)$.
Continuity at $\rho = 1$ follows by L'H\^opital's rule.
Full details are in Appendix~\ref{app:proofB}.
\end{proof}

\begin{figure}[ht]
\centering
\begin{knitrout}
\definecolor{shadecolor}{rgb}{0.969, 0.969, 0.969}\color{fgcolor}
\includegraphics[width=0.84\textwidth]{figures/knitr-bridge-cutoff-figure-1} 
\end{knitrout}
\caption{The bridge cutoff as a function of preferences and the
long-asset return. The numerical value of $p^*$ is calibration
dependent; the theorem establishes the threshold form rather than a
universal cutoff.}
\label{fig:bridge_cutoff}
\end{figure}

\subsection{Interpretation}

The theorem translates the canonical decision into a form suitable
for the next model. When the insurance premium is small, the
withdrawal decision can be expressed by asking whether estimated
recovery exceeds a fixed threshold. This is the formal handoff from
the DD expected-utility problem to the recovery comparison used in
the network model; it is not a separate behavioral postulate. The
theorem's critical implication is structural: under its binary-payment
conditions, an expected-utility calculation reduces to a $p^*$ rule.
It does not imply that one numerical value of $p^*$ is universal.
The cutoff varies with the utility and payoff calibration (and, under
other normalizations, may vary with claim size).
We use \emph{recovery-probability cutoff} for this theorem-defined
family and reserve \emph{calibrated bridge threshold} for the
numerical value 0.80 used in the network experiment.

In the limiting case of a complete observation network---where every
agent observes every other agent's withdrawal decision---the local
survival estimate equals the true value $P(\text{survival} \mid D)$,
and the bridge theorem applies directly though approximately: under
the theorem's maintained primitives, the individually rational
decision rule is to withdraw if and only if estimated survival falls
below $p^*(\rho, R)$. The complete-graph case thus provides a
microfoundation for the threshold rule. For sparse networks ($k \ll K$),
agents form beliefs from a strict subset of the depositor pool; the
threshold rule remains individually rational given those local beliefs,
but the departure from the complete-graph benchmark is an information
friction---agents act on incomplete signals---not a coordination failure
among fully informed agents.

The convergence is quantitatively first-order:

\begin{corollary}[First-order approximation]
\label{cor:approx}
Under the conditions of Theorem~\ref{thm:bridge}, convergence is
linear in $\varepsilon$:
\begin{equation}
  \Phi_{\rm shift}(\varepsilon)
  \;=\; p^*(\rho,\,R)
        \;+\; M(\rho,\,R,\,\pi)\cdot\varepsilon
        \;+\; O(\varepsilon^2)
  \label{eq:approx}
\end{equation}
where $M(\rho,R,\pi) = d\Phi_{\rm shift}/d\varepsilon\big|_{\varepsilon=0} > 0$.
For $\rho \neq 1$, set $\alpha = 1-\rho$, $A = (1+R)^\alpha - 1$,
$B = 2^\alpha - 1$, and $\mu = \pi R/(1-\pi)$.
In this notation $p^*(\rho,R) = B/A$ (from equation~\eqref{eq:pstar}),
and:
\begin{equation}
  M(\rho,R,\pi) \;=\;
  \frac{\alpha\bigl[2^{\alpha-1}A \;+\; \mu\,(1+R)^{\alpha-1}B\bigr]}
       {A^2},
  \qquad \rho \neq 1.
  \label{eq:Mformula}
\end{equation}
For $\rho = 1$, $\alpha = 0$ makes $A$ and $B$ degenerate;
differentiating the log-utility form directly gives:
\begin{equation}
  M(1,R,\pi) \;=\;
  \frac{\dfrac{\log(1+R)}{2} \;+\; \mu\,\dfrac{\log 2}{1+R}}
       {\log(1+R)^2},
  \qquad \mu = \frac{\pi R}{1-\pi}.
  \label{eq:Mformula_log}
\end{equation}
Since $M > 0$, $\Phi_{\rm shift}(\varepsilon) > p^*(\rho,R)$ for all
$\varepsilon > 0$: under DD's mechanism the run threshold strictly
exceeds the limit for any positive premium.
\end{corollary}

The sign $M > 0$ is the same $\alpha^2$ argument as
Lemma~\ref{lem:phi} (see Appendix~\ref{app:proofB}): since
$M > 0$, $\Phi_{\rm shift}(\varepsilon) > p^*(\rho,R)$ for all
$\varepsilon > 0$, confirming that the withdrawal threshold strictly
exceeds its zero-premium limit for any positive premium.

Theorem~\ref{thm:main} gives a conditional direction, while
Corollary~\ref{cor:approx} describes the local change in the cutoff;
neither alone supplies the magnitude of failure. The fixed-contract
simulation supplies that distributional object. It shows an aggregate
inward movement of the failure boundary, with the local reversals
reported in Appendix~\ref{app:ddsweep}.

The paired bridge experiment addresses a different question: the
withdraw/stay classification near the zero-premium limit. Its two
partial-payment bounds agree in 153 of 180
scenarios. The allocation panel separately measures participation.

\paragraph{Generalization to risky withdrawal.} Theorem~\ref{thm:bridge}
assumes immediate withdrawal succeeds with certainty: the withdrawing
agent receives $c_1$ for sure. In the network model below, immediate
withdrawal may itself be risky---an agent who withdraws may be late in
the queue and face depleted reserves. Let $p_W$ denote the probability
of full recovery under immediate withdrawal and $p_S$ the probability
of full recovery under waiting. The utility-consistent comparison is:
%
\begin{equation}
  p_W\,u(c_1) + (1-p_W)\,u(0)
  \;\geq\;
  p_S\,u(c_2) + (1-p_S)\,u(0)
  \label{eq:generalbridge}
\end{equation}
%
Rearranging gives the generalized withdrawal condition:
%
\begin{equation}
  p_S \;\leq\; p_W\,\Phi(c_1,\,c_2,\,\rho)
  \label{eq:generalcond}
\end{equation}
%
As $\varepsilon \to 0$ this converges to:
%
\begin{equation}
  p_S \;\leq\; p_W\,p^*(\rho,\,R)
  \label{eq:generalcondlimit}
\end{equation}
%
Theorem~\ref{thm:bridge}'s threshold rule $p_S \leq p^*(\rho,R)$
is the special case $p_W = 1$.
This case has a particularly transparent interpretation: when
withdrawing pays in full with certainty, the depositor withdraws
exactly when the probability of full recovery from staying is no
greater than the utility-derived cutoff $p^*$. At the other boundary,
if $p_S=1$, equation~\eqref{eq:generalcondlimit} cannot hold when
$R>1$: because $p_W\leq1$ and $p^*(\rho,R)<1$, we have
$p_Wp^*(\rho,R)<1=p_S$. Thus certain full recovery from staying
implies staying. These boundary cases also explain the certainty
condition imposed in the network implementation.

The network simulations below evaluate this implication rather than
assuming it. They compare the relative-safety rule $p_W>p_S$, the
bridge threshold in equation~\eqref{eq:generalcondlimit}, and an
explicit expected-utility calculation based on paired withdraw-now
and stay payments. Common scenario and model seeds make the
comparison within-realization. This design separates the
information-contagion mechanism---local observations alter perceived
recovery probabilities---from the particular rule mapping those
probabilities into an action.

%%============================================================
\section{Bridge-Theorem Simulation}
\label{sec:bridgesimulation}
%%============================================================

The finite-agent bridge evaluates the limiting withdraw-versus-stay
comparison under sequential service. It contains 180 scenarios and
10{,}000 paired trials per scenario. In every trial, withdrawing now
and staying face the same draw of other agents' withdrawals. Because
sequential service can generate partial payments, the analysis reports
two deliberately extreme classifications: a withdrawal-favoring bound
and a staying-favoring bound.

The two bounds agree in 153 of 180 scenarios.
Both prescribe withdrawal in 75 scenarios and
staying in 78; only 27
scenarios depend on how partial payments are classified. This
calculation uses homogeneous finite-agent DD primitives and paired
counterfactual payments rather than assuming the network decision
rule. It therefore evaluates whether the theorem's limiting decision
object remains stable in a finite sequential-service queue.

\begin{figure}[ht]
\centering
\begin{knitrout}
\definecolor{shadecolor}{rgb}{0.969, 0.969, 0.969}\color{fgcolor}
\includegraphics[width=0.9\textwidth]{figures/knitr-bridge-classification-figure-1} 
\end{knitrout}
\caption{Finite-agent bridge classifications. Points report whether
the paired withdraw-now and stay comparison prescribes withdrawal or
staying under both extreme treatments of partial payment. Gray points
identify scenarios whose classification changes with the bound.}
\label{fig:bridge_classification}
\end{figure}

\begin{table}[ht]
\centering
\caption{Finite-agent bridge decisions under partial-payment bounds}
\label{tab:bridge_bounds}

\begin{tabular}{lr}
\toprule
Classification & Scenarios\\
\midrule
Both bounds: withdraw & 75\\
Both bounds: stay & 78\\
Bound-sensitive & 27\\
\bottomrule
\end{tabular}


\end{table}

%%============================================================
\section{Network-Contagion Model}
\label{sec:replacement}
%%============================================================

\subsection{Motivation}

Conditional on its assumed survival map, Theorem~\ref{thm:main}
shows that the withdrawal threshold rises and failure probability
weakly increases with the insurance premium $\varepsilon$, with
strictness under the stated positive-probability crossing condition.
Theorem~\ref{thm:bridge} shows that as $\varepsilon \to 0$, rational
CRRA agents reduce to probability threshold comparators. Real deposit
contracts have $\varepsilon \approx 0$. The incremental above-par
component is therefore absent or small, while the recovery-priority
incentive created by sequential service remains.

Sequential service is common to the canonical and network models;
the marginal payoff and information mechanisms are not. It is
necessary for the recovery-priority mechanism studied here, though
we do not claim it is necessary for every possible model of a bank
run. In the canonical DD contract, the above-par early payment changes
the patient depositor's marginal payoff comparison. Under par
redemption, sequential service continues to make recovery depend on
priority, while local observations additionally change beliefs about
future withdrawals.

What, then, helps determine marginal withdrawal under par redemption?
We propose that some real bank runs are
information cascades on depositor networks. Agents observe local
withdrawal behavior, treat it as a signal about global conditions, and
act when their estimated probability of loss crosses the threshold $p^*$
derived in Theorem~\ref{thm:bridge}. This differs from
the strategic multiplicity in Diamond--Dybvig without eliminating
strategic interaction. In the DD model, agents are
fully informed about fundamentals; multiple equilibria arise because each
agent's optimal action depends on what others will do, independently of
information. In the network model, the source of runs is different: at
any given activation step, conditional on the current deposit pool $D(t)$
and the agent's local signal $\hat{w}_i(t)$, the withdrawal decision is
individually rational given the agent's beliefs---the agent withdraws
because their estimated loss probability exceeds $p^*$, rather than
because of an exogenous sunspot. The process nevertheless retains
strategic complementarity and is path-dependent: $D(t)$ is endogenous
to all prior withdrawals, and the Monte Carlo survival estimate
incorporates beliefs about future
withdrawals drawn from the truncated geometric prior. Earlier withdrawals
therefore feed forward into later agents' signals and probability
estimates. They make later withdrawal more attractive by reducing the
available pool and changing local information. Runs arise because
agents are \emph{not} fully informed: each agent observes only a
neighborhood of size $k < K$, and local withdrawals can push the
estimated survival below $p^*$ even when the bank is fundamentally
sound. The
mechanism is incomplete information propagating through a network.
\subsection{Setup}
\label{sec:setup}

There are $K$ agents on a fixed undirected network $G = (V, E)$.
Agent $i$ has deposit $d_i$ drawn independently from a log-normal
distribution $F = \mathrm{LogNormal}(\mu, \sigma)$. A deposit
insurance threshold $\bar{d}$ fully insures agents with $d_i \leq
\bar{d}$; in the simulation, $\bar{d}$ is set to the $q$-th quantile
of $F$ for parameter $q \in [0,1)$. Insured agents have no incentive
to withdraw regardless of network signals and are excluded from the
cascade dynamics. All statements below concern uninsured agents
($d_i > \bar{d}$).

The bank holds reserves equal to a fraction $r$ of total deposits.
Agent $i$ observes the local withdrawal rate in their neighborhood
$\mathcal{N}(i)$:
%
\begin{equation}
  \hat{w}_i(t) \;=\;
  \frac{\bigl|\{j \in \mathcal{N}(i) : j \text{ has withdrawn by }
  t\}\bigr|}{|\mathcal{N}(i)|}
  \label{eq:localrate}
\end{equation}
%
Agent $i$ treats $\hat{w}_i(t)$ as representative of the global
withdrawal rate and forms a belief about total withdrawals. Formally,
the agent treats the total number of withdrawals as drawn from a
Geometric distribution with parameter $p_0 = 0.1$, truncated from
below at the locally implied count:
%
\begin{equation}
  \tau_i(t) \;=\; \mathrm{round}\!\left(K \cdot \hat{w}_i(t)\right)
  \label{eq:tau}
\end{equation}
%
\begin{equation}
  W_i \;\sim\; \mathrm{Geometric}(p_0)\;\big|\;
  W_i \geq \tau_i(t)
  \label{eq:belief}
\end{equation}
%
The Geometric is the unique discrete memoryless distribution. Of course,
the truncation introduces some memory but the result still holds
approximately.
The agent uses $\tau_i(t)$ as a lower bound on total withdrawals,
conditioning the fixed-parameter geometric on the observed count;
this is not Bayesian updating over $p_0$, but a restriction of the
support given what the agent knows.

\paragraph{Belief-model qualification.} The truncated Geometric is a
deliberately parsimonious baseline chosen for analytical tractability
and computational transparency. It is neither a uniquely rational
posterior nor an empirically estimated depositor belief process.
First, the memoryless
property means the conditional distribution of $W_i$ depends only on
the current lower bound $\tau_i(t)$, not on the history of how that
bound was reached; this is appropriate when agents have no memory of
prior activation states. Second, the single parameter $p_0$ has a direct
interpretation as the prior probability that any given agent
withdraws, and the uniform choice $p_0 = 0.1$ for both the
exogenous shock process and the belief prior imposes internal
consistency between the two stochastic objects in the model. Third,
the closed-form survival probability $p_S(x)$ derived in
Appendix~\ref{app:abm} follows directly from the geometric's
closed-form tail, which keeps the agent decision problem
analytically tractable.

Full Bayesian updating over the withdrawal probability $p_0$ would
require agents to maintain and propagate a posterior distribution
over $p_0$ at each activation step across $K = 1{,}000$ agents and
would be substantially more demanding in the present implementation.
The truncated Geometric is used as the tractable baseline.

\paragraph{Information-structure qualification.} Agents use local
observations as a proxy for global conditions because global depositor
network information is unavailable. Conditional on the maintained
belief model, this is an internally coherent response to censored
information, not an assertion that the proxy is empirically optimal.
SVB depositors observed
withdrawals within their venture-capital and startup networks; their
local sample was the best available signal, even though it was a
non-representative sample of the full depositor base.

The withdrawal decision is implemented as a pairwise safety comparison.
The agent simulates two scenarios by Monte Carlo, drawing from
belief~\eqref{eq:belief}: (i)~withdraw now and take position in the
current queue; (ii)~stay and wait for the process to settle. In each
draw, the agent records whether they recover their full deposit. Let:
%
\begin{align}
  p_W &\;=\; P(\text{recover } d_i \mid \text{withdraw now},\,
              \hat{w}_i(t)) \notag \\
  p_S &\;=\; P(\text{recover } d_i \mid \text{stay},\,
              \hat{w}_i(t)) \notag
\end{align}
%
The simulation evaluates three mappings from these paired
counterfactuals into withdrawal. The comparative rule withdraws if
$p_W>p_S$. The bridge rule uses the generalized condition
(equation~\eqref{eq:generalcond}):
%
\[
  p_S \;\leq\; p_W\,p^*(\rho,\,R).
\]
%
The explicit-utility rule instead compares the expected shifted-CRRA
utility of the full paired payment distributions, retaining zero,
partial, and full recovery rather than reducing each draw to a binary
indicator. All three rules first impose a certainty condition: the
agent stays when full recovery from waiting is certain.

The comparative rule remains useful as a transparent benchmark.
When the loss of principal dominates the interest differential in
utility terms---which holds for any strictly concave $u$ when
$c_2-c_1$ is small relative to $c_1$---the ratio
%
\[
  \frac{u(c_2)-u(0)}{u(c_1)-u(0)} \;\approx\; 1
\]
%
and the exact rule and the comparative approximation may nearly
coincide. Section~\ref{sec:results} reports where this robustness
holds and, equally importantly, the selected boundary environments
where it does not.

\subsection{Activation Process}

Exogenous withdrawals occur at the start of each simulation run.
The number of exogenously withdrawing agents is drawn from a
$\mathrm{Geometric}(p_0)$ distribution truncated to $\{0,\ldots,K\}$,
and the withdrawing agents are selected uniformly at random.

Following the exogenous phase, banking agents are processed in a
uniformly random sequential order each period. Each activated agent
computes $\hat{w}_i(t)$ from the current neighborhood state and applies
the comparison rule. Periods repeat until a full pass over all banking
agents produces no new withdrawals---a fixed point of the cascade
dynamics.

\subsection{Cascade Endpoints and an Analytical Threshold Benchmark}

\begin{definition}[Terminal cascade state]
\label{def:equilibrium}
A withdrawal configuration $\mathbf{w}^{\dagger}\in\{0,1\}^K$ is a
\textit{terminal cascade state} if a full sequential pass over all
agents who have not withdrawn produces no additional withdrawal.
\end{definition}

This definition matches the simulation's halting criterion. Because
withdrawal is irreversible in the implementation, a terminal cascade
state is an absorbing endpoint of the one-way activation process; it
need not be a Nash equilibrium in which previously withdrawn agents
are allowed to reverse their actions.

For comparison, consider an idealized deterministic threshold map
$T^0$ that evaluates loss probabilities exactly, fixes insurance and
other bank-state components, and withdraws when a monotone local-loss
signal exceeds a fixed cutoff. Conditional on these restrictions,
more neighboring withdrawals weakly increase every affected signal,
so $T^0$ is isotone. Tarski's theorem then gives minimal and maximal
fixed points for this analytical benchmark. Under the additional
conditions that the exogenous-only signal lies below the cutoff and
the all-neighbors-withdrawn signal lies above it, these fixed points
include a low-withdrawal and a high-withdrawal rest point.

This lattice statement does not extend automatically to the realized
simulation. Its three decision rules use finite Monte Carlo
counterfactuals, partial payments, reserve depletion, and in some
specifications changing insurance coverage. Sampling error can also
flip a realized classification. We therefore use the lattice result
only to describe the exact deterministic threshold benchmark, not to
claim isotonicity or generic multiplicity for every simulated state
or decision rule. The reported computational results concern the
distribution of terminal cascade states generated by the implemented
sequential process.

The simulation uses Newman--Watts--Strogatz (NWS) networks
\citep{newman1999renormalization}, which add shortcut edges to a ring
lattice with probability $p_{\rm ws}$ per pair, keeping the lattice
edges intact. Unlike standard Watts--Strogatz rewiring, NWS networks
remain connected and preserve local clustering while introducing
long-range links. The parameters $k$ (ring lattice degree) and
$p_{\rm ws}$ control the clustering--path-length tradeoff and are
varied in the comparative statics of Section~\ref{sec:results}.

\paragraph{Activation order and cascade paths.}
Which terminal cascade state is reached depends on the order in which
agents are activated. High-withdrawal endpoints are more likely when
high-degree agents activate early in the sequence: an
early withdrawal by a high-degree agent shifts the estimated loss
probability for many neighbors simultaneously, potentially producing
a larger cascade. This activation-order dependence is
structural within the model and corresponds to a testable account of
how observable early movers may trigger cascades among their network
neighbors. Unlike DD, where equilibrium selection may be
represented by an exogenous sunspot, path selection here is generated
by the activation order, local sample size, and shortcut density. The simulation
draws activation order independently for each run; the distribution of
outcomes across draws defines the run probability reported in
Section~\ref{sec:results}.

%%============================================================
\section{Network Results}
\label{sec:results}
%%============================================================

\subsection{Simulation Design}

The network evidence uses three rules on matched scenarios:
(i) a comparative rule that withdraws when full recovery is more
likely from withdrawing than from staying; (ii) the bridge threshold;
and (iii) explicit expected shifted-CRRA utility computed from paired
payment draws. The bridge rule is a reduced-form member of the
threshold family justified by Theorem~\ref{thm:bridge}; its numerical
cutoff of 0.80 is the binary-recovery, risk-neutral calibration with
par withdrawal and gross late payoff 1.25:
\[
  p^*=\frac{1}{1.25}=0.80.
\]
It is not the universal value of the theorem's cutoff family. For
example, shifted log utility with $a=1$ and the same two successful
payoffs gives $\log(2)/\log(2.25)\simeq0.855$. The explicit-utility
rule retains the
full paired payment distribution and therefore checks whether the
network findings depend materially on that reduced-form calibration.
All three rules impose the certainty condition that an agent stays
when full recovery from waiting is certain, as the generalized bridge
comparison requires when $R>1$. The primary
production evidence is a 39{,}600-run, outcome-independent,
prespecified monotonicity sweep containing six deliberately varied
deposit/topology structures, eleven settings per structure, three
decision rules, and 200 replications. Stable common seeds match the
three rules within every structural cell and replication. An adaptive
boundary stress test is reported separately in
Appendix~\ref{app:adaptive}.

The validation, monotonicity, and adaptive exercises answer different
questions. The validation uses a compact common-scenario grid; the
prespecified monotonicity design uses balanced one-factor paths in
reserves, fixed insurance, and shock size across six deliberately
varied structures; and the adaptive stress test selects high-entropy,
disagreement-prone cells. Their agreement rates are therefore
conditional on their respective designs and are not estimates from a
common representative parameter distribution. In particular, the
monotonicity design was frozen without using outcomes from its cells;
it should not be described as an outcome-selected sample of transition
boundaries.

\subsection{Local Information, Sequential Service, and Decision Rules}

The network results establish the sufficiency claim directly.
Depositors observe only realized withdrawals in their neighborhoods,
do not observe their global queue positions, and form beliefs about
unobserved withdrawals. Payments obey sequential service. Under this
information structure, every decision rule produces both endogenous
withdrawals and bank failures in the prespecified sweep.
Table~\ref{tab:network_rules} reports the unconditional rates. The
failure rates differ across rules, but endogenous contagion is not
created by a particular shortcut to the decision problem: it occurs
under the comparative, bridge-threshold, and explicit-utility rules.

\begin{table}[ht]
\centering
\small
\caption{Network outcomes by decision rule: prespecified sweep}
\label{tab:network_rules}

\begin{tabular}{lrrr}
\toprule
Decision rule & Failure rate & Pr(contagion) & Mean endogenous\\
\midrule
Explicit utility & 0.111 & 0.376 & 22.2\\
Comparative recovery & 0.068 & 0.450 & 36.9\\
Bridge threshold & 0.149 & 0.403 & 25.2\\
\bottomrule
\end{tabular}


\smallskip

\small\textit{Note}: 13{,}200 runs per decision rule, matched across
66 structural cells and 200 replications. In the network experiment,
``failure'' denotes vault exhaustion under sequential service. In the
DD experiment, ``default'' instead denotes an underpaid contractual
claim; exact exhaustion after full payment is settlement.
\end{table}

The closest comparison for the paper's robustness claim is the
bridge threshold against explicit utility. On the 600 matched
validation scenarios, the rules agree on the bank-failure outcome in
99.5\% of cases. Their failure
rates are 8.2\% and
8.7\%, a difference of only
0.5\%. In the broader
prespecified monotonicity sweep, outcome agreement remains
96.1\% and the aggregate
failure-rate difference is
3.9\%.

High classification agreement does not imply identical aggregate
failure rates. In a path-dependent cascade, disagreements concentrated
near a transition can be systemically consequential: changing an early
withdrawal changes both the remaining reserve pool and the signals
observed by agents activated later. A small share of paired runs can
therefore account for a meaningful directional difference in aggregate
failure rates.

This robustness has a clear boundary. The two specifications are not
numerically identical: one is a theorem-motivated reduced-form
threshold rule and the other computes expected utility from the full
paired payment distribution. The adaptive stress test in
Appendix~\ref{app:adaptive} deliberately concentrates uncertain and
model-disagreement cells and shows that the exact rule can matter near
contested decision boundaries. The supported conclusion is therefore
narrower than decision-rule irrelevance: local information and
sequential service generate runs under each microfoundation, while
the bridge threshold and explicit utility give nearly the same
failure classification in ordinary validation and most prespecified
settings.

These are mechanism results conditional on the specified beliefs,
links, activation, and payments. They neither calibrate actual
depositor networks nor validate the truncated-Geometric belief law.
Section~\ref{sec:predictions} states the data needed for empirical
identification. Policy welfare is outside the simulation exercise.

\begin{figure}[ht]
\centering
\begin{knitrout}
\definecolor{shadecolor}{rgb}{0.969, 0.969, 0.969}\color{fgcolor}
\includegraphics[width=0.88\textwidth]{figures/knitr-network-validation-figure-1} 
\end{knitrout}
\caption{Common-scenario network validation. Each panel holds the
initial exogenous withdrawal count fixed. Failure rates fall sharply
with the reserve ratio under all three local-information decision
rules. The bridge threshold and explicit-utility curves nearly
coincide.}
\label{fig:network_validation_new}
\end{figure}

Figure~\ref{fig:network_validation_new} also clarifies the role of
fundamentals. Local information is sufficient to propagate runs, but
does not make reserves irrelevant: failure declines sharply as the
reserve ratio rises. The network mechanism is therefore neither a
claim that every local signal causes a run nor a claim that topology
dominates bank resources. It is an existence and robustness result:
given sequential service, locally observed withdrawals can generate
self-reinforcing endogenous withdrawals without common knowledge of
fundamentals or global queue position.

\begin{figure}[ht]
\centering
\begin{knitrout}
\definecolor{shadecolor}{rgb}{0.969, 0.969, 0.969}\color{fgcolor}
\includegraphics[width=0.92\textwidth]{figures/knitr-network-comparative-statics-figure-1} 
\end{knitrout}
\caption{Prespecified network comparative statics. Each panel
averages the balanced one-factor path over six deliberately varied
deposit and topology structures and 200 replications per structural
cell. The panels are design-specific comparative statics, not
population-weighted estimates.}
\label{fig:network_comparative_statics}
\end{figure}

% Legacy empirical block 1 archived in archive/legacy_empirical_blocks_20260724.tex.

%%============================================================
\section{Contrast and Policy Implications}
\label{sec:policy}
%%============================================================

\subsection{Two Models of the Same Phenomenon}

Both models produce coordination failure, support two equilibria,
and generate bank runs in which individually rational behavior
produces collectively destructive outcomes. The distinction is
mechanism. Table~\ref{tab:contrast} summarizes the structural
differences.

\begin{table}[ht]
\centering
\caption{Structural comparison: Diamond--Dybvig, global games,
  and network-contagion model}
\label{tab:contrast}
\small
\begin{tabular}{lp{2.4cm}p{2.4cm}p{2.6cm}}
\hline\hline
 & DD Model & Global Games & Network-Contagion \\
\hline
Run mechanism        & Sequential service plus premium & Fundamental threshold       & Local-information cascade \\
Key marginal driver  & Premium $\varepsilon$     & Signal noise $\sigma$       & Recovery beliefs \\
Equilibrium          & Sunspot (arbitrary)       & Unique                      & Activation order (structured) \\
Decision rule        & CRRA utility max          & Cutoff on private signal    & Behavioral approx.\ $p_W\!>\!p_S$ \\
Contract environment & Canonical $c_1>1$ contract & Standard deposits          & Par-value deposits \\
Run interpretation   & Calibrated CRRA           & Calibrated                  & Comparative statics \\
Parameters observable& Contract specific         & Partial ($\sigma$ unobservable) & Partial ($k$, $r$, $q$) \\
\hline\hline
\end{tabular}
\smallskip

\small\textit{Note}: Canonical DD admits fundamental and panic
equilibria. Under the conditions stated in
Section~\ref{sec:replacement}, the deterministic network-threshold
benchmark admits low- and high-withdrawal rest points; the simulation
reports irreversible terminal cascade states. Global games yields a
unique threshold equilibrium. The rows contrast mechanisms,
information environments, and observability.
\end{table}

The three models are not competing accounts of whether runs can
happen---all three confirm they can. They are competing accounts of
why runs happen and which structural features determine
their frequency. That distinction matters for policy: the correct
intervention target depends entirely on which mechanism is
operative.

Global games \citep{goldstein2005demand} and the network-contagion
model address the multiplicity problem in DD from opposite directions.
Global games eliminates multiplicity by introducing heterogeneous
private signals about bank fundamentals: noise in those signals pins
down a unique threshold equilibrium. The network simulation instead
uses a structural path mechanism---activation order, neighborhood
degree, and shortcut density---that is, in principle, observable. These are
complementary contributions operating in different information
environments, as discussed in the Introduction. They also generate
distinct empirical predictions: in global games, run probability
is fully determined by fundamentals and signal precision, so two
banks with identical fundamentals face identical run risk regardless
of depositor network structure. In the present NWS specification,
neighborhood degree $k$ and shortcut density predict run probability
conditional on fundamentals. This is testable using depositor-link
and withdrawal-timing data and is not a standard prediction of the
baseline global-games framework.

\subsection{Mechanism Mismatch in the Marginal Withdrawal Condition}

The canonical DD optimal contract and an ordinary demand deposit do
not give patient depositors the same marginal withdrawal problem. In
the canonical contract, $\varepsilon=c_1-1>0$ enters the comparison
directly and also lowers continuation consumption through the
resource constraint. Conditional on the imposed survival map,
Theorem~\ref{thm:main} establishes that the withdrawal threshold rises
and failure probability weakly increases with this premium. An
ordinary demand deposit instead redeems
at par. It therefore removes the incremental above-par channel but
retains sequential service, recovery priority, and strategic
complementarity.

Real deposit contracts do not embed $\varepsilon > 0$. FDIC
insurance pays par. Time deposits carry early withdrawal penalties
that move $c_1$ below par. The instruments that most closely
approximate DD's contract---money market funds, overnight repo,
certain structured investment vehicles---are precisely those that
experienced destabilizing runs in 2008 and 2020.\footnote{Three
  instrument types can be distinguished for applying these
  mechanisms. Type~1 ($\varepsilon > 0$): instruments with
  contractual above-par early redemption; DD's option-exercise
  mechanism is fully operative.
  Type~2 ($\varepsilon = 0$, sequential service): standard
  demand deposits; the above-par component is absent, while
  recovery-priority and information mechanisms remain operative.
  Type~3 (repo, MMF): par-redemption instruments where runs
  take the form of rollover refusal or gate imposition rather than
  above-par redemption \citep{gorton2012securitized}; closer to DD
  than retail deposits---haircut increases approximate
  $\varepsilon \to 0^{-}$---but the run is triggered by counterparty
  refusal rather than depositor option exercise, placing them in
  mixed-mechanism territory.}
This pattern motivates separating the incremental premium channel
from the recovery-priority logic shared by wholesale and retail
claims.

The quantitative record supports this scope delineation. Standard
retail certificates of deposit carry penalties of 90 to 180 days'
simple interest for maturities of one to two years
\citep{reg_d, reg_dd},
placing $\varepsilon$ firmly negative for time deposits and exactly
zero for demand deposits. FDIC deposit insurance pays par by statute
\citep{fdia_1821f}, removing any above-par premium for insured
balances. The repo run documented by \citet{gorton2012securitized}
involved haircut increases rather than above-par redemption: the
mechanism was creditor refusal to roll at par, not option exercise.
Even setting $\varepsilon = 0.02$---above any documented positive
premium in retail banking---Corollary~\ref{cor:approx} implies the
withdrawal threshold rises only modestly above its zero-premium
limit. The paired bridge results show that the limiting decision is
robust to the treatment of partial recovery in 153 of 180
finite-agent scenarios.

DD is a foundational theory of coordination failure under sequential
service and a specific liquidity-insurance contract. The
contract-structure mismatch limits direct transfer of its canonical
marginal condition to ordinary retail deposits: standard deposits do
not provide the above-par payoff. The network model retains what is
shared---sequential service and recovery priority---and adds a local
information mechanism suited to par-value deposits. It is therefore
a refinement of the marginal environment, not a rejection of DD's
broader fragility logic.

\subsection{Policy Implications}

DD's canonical policy prescription is deposit insurance:
government guarantees make early withdrawal unprofitable by
eliminating the downside risk that motivates option exercise. This
prescription is sound within DD's mechanism and remains relevant to
retail banking through the shared recovery-priority problem. What
does not transfer unchanged is the incremental above-par incentive
embedded in the canonical optimal contract.

Two critiques should be kept separate here. The first is a
critique of DD's mechanism: the option-exercise channel
is absent from ordinary deposit contracts, so it cannot be the full
account of marginal withdrawal behavior in retail banking. The
second is a separate question about whether deposit
insurance succeeds in our network-contagion model, where
the option channel is absent by construction. These questions have
different answers. In Section~\ref{sec:setup}, insured agents are
excluded from the cascade dynamics entirely: because $p_W = p_S =
1$ for an insured agent regardless of network signals, the
withdrawal condition $p_W > p_S$ is never satisfied, and insured
agents never run. Deposit insurance therefore reduces run
probability in our model through a distinct channel---removing
agents from the cascade pool rather than eliminating an option
payoff. The policy works, but the mechanism is participation
suppression rather than option-value elimination, and the two
explanations carry different empirical implications. The DD account
predicts that run probability is increasing in the size of the
above-par premium; our account predicts that it is decreasing in
the share of deposits covered, which is directly reflected in the
reserve ratio and insurance quantile sweeps of
Section~\ref{sec:results}. The SVB episode is consistent with our
account: uninsured balances above the FDIC threshold faced genuine
loss risk, deposit insurance was not credibly available for those
balances, and the cascade propagated through the uninsured
depositor network.

The network-contagion model points to a different and complementary
set of intervention levers.

\paragraph{Neighborhood degree and shortcut density.}
In the reported NWS comparison, run probability decreases with
neighborhood degree $k$. The operative margin is local sample size,
not a claim that every denser topology is safer. At adequate
reserves ($r = 0.4$), the run rate falls from 78.5\% at $k = 6$ to
43.6\% at $k = 50$. The mechanism is sample quality: agents in
larger neighborhoods face less sampling variance and are less likely
to overestimate the global withdrawal rate in this specification.
Larger samples need not be unbiased when network position and
withdrawal behavior are correlated. Shortcut density separately
changes clustering and path length, so it should not be conflated
with degree.

The model-implied hypothesis runs in the opposite direction:
institutions
serving highly concentrated depositor bases face structural cascade
risk unrelated to contract terms. SVB's depositor base, concentrated
among venture-capital firms and their portfolio companies with dense
mutual-observation networks, exemplifies this risk
\citep{cookson2023social}. Supervisory attention to depositor
concentration---particularly for uninsured balances above the FDIC
threshold---arises directly in the network model rather than the
canonical DD premium channel.

\paragraph{Information architecture.}
The cascade mechanism operates through the observability of
withdrawal decisions. Agent $i$ withdraws when their local signal
$\hat{w}_i(t)$ exceeds the threshold $p^*$; that signal is
informative only because withdrawals are observable to network
neighbors. This makes the cascade sensitive to policies that reduce
the informativeness of early-mover signals.

Consider a \textit{sequenced redemption window}: agents are
partitioned into batches of size $B$ and activated simultaneously
within each batch, so that agents in a given round observe only
withdrawals from previous rounds, not from neighbors in the same
round. Batching suppresses the cascade through two channels.
First, within-batch contagion is eliminated: a single early
withdrawal cannot update the signal of a neighbor activated in the
same round, so the path-dependent amplification that drives cascade
dynamics is broken at the within-batch level. Second, cross-batch
signals convey aggregate rather than individual withdrawal rates;
the inference problem shifts from ``my neighbor withdrew'' (high
signal variance, easily tipped by idiosyncratic shocks) to
``fraction $x$ of the previous cohort withdrew'' (lower variance,
more reflective of fundamentals). In the limit where $B = K$ and
all agents are activated simultaneously, no agent observes any
endogenous withdrawal before deciding; the cascade cannot propagate
and run probability collapses to the exogenous shock rate $p_0$.
This mechanism has direct counterparts in the existing regulatory
toolkit. SEC Rule~22e-3 under the Investment Company Act authorizes
money market fund boards to suspend redemptions when weekly liquid
assets fall below 30\% \citep{sec_rule_22e3}, and the 2016
institutional MMF reforms extended the toolkit to include floating
net asset value requirements and discretionary liquidity fees
\citep{sec_mmf_reform_2016}. Redemption gates
implement the batching logic above: by suspending observable
redemptions, they break the signal that propagates the cascade.
\citet{ennis2009institutions} caution that a discretionary
suspension power is itself a source of strategic risk: if
depositors cannot be sure the regulator will refrain from
intervening, the anticipated intervention can change withdrawal
incentives ex ante, so the instrument's effect on run risk is not
simply the direct effect of suspending redemptions in isolation.
Floating NAV removes the at-par redemption feature that brings
wholesale instruments closest to the DD contract structure, moving
the effective $\varepsilon$ toward zero. Whether these instruments
improve welfare or are optimally calibrated is a distinct question;
the mechanism that
motivates them---degrading the cascade-relevant signal---is the
same one identified by our model.

\paragraph{Discretionary triggers versus universal batching.}
The 2023 MMF reforms illustrate why the batching mechanism above
specifically requires that suspension be \textit{universal and
pre-announced} rather than discretionary and threshold-triggered.
The SEC eliminated the redemption gate entirely and severed the tie
between discretionary liquidity fees and the 30\% weekly-liquid-assets
threshold, replacing it for institutional prime and institutional
tax-exempt funds with a mandatory liquidity fee triggered by daily net
redemptions exceeding 5\% of net assets \citep{sec_mmf_reform_2023}.
The adopting release attributes this reversal directly to the March
2020 experience: even though no fund actually imposed a fee or gate,
``the possibility of their imposition after crossing the publicly
disclosed 30\% weekly liquid asset threshold appears to have
contributed to investors' incentives to redeem''
\citep{sec_mmf_reform_2023}. In our terms, the 30\% threshold was
itself an observable signal---proximity to the trigger---that
depositors could watch and act on preemptively, which is exactly the
kind of endogenous, state-dependent signal that seeds a cascade in our
model. This is a different channel from the one our batching proposal
targets. A discretionary threshold gate adds a new observable state
variable (distance to the trigger) on top of whatever withdrawal
information already propagates through the network; agents who can
monitor that state variable acquire an additional coordination device.
Universal, pre-announced, simultaneous batching does not have this
feature: because every agent is suspended on the same fixed schedule
regardless of the fund's or bank's realized liquidity position, there
is no threshold-proximity variable to observe or race against. The
2023 reversal is therefore consistent with our mechanism rather than
in tension with it: it removed a discretionary, state-triggered
suspension power precisely because that power added a signal, and our
proposal is deliberately designed not to add one.

These instruments target an information-propagation margin that is
not explicit in the canonical DD framework. DD emphasizes the payoff
structure of the deposit contract: deposit insurance eliminates
recovery risk and thereby weakens the incentive to withdraw early. A
redemption gate can additionally change which withdrawals become
observable before later agents decide. The network model therefore
motivates a complementary family of policy hypotheses with
counterparts in the regulatory toolkit for both retail deposits and
wholesale funding instruments. Evaluating a particular intervention
requires additional analysis of anticipation, welfare, commitment,
and the possibility that delayed disclosure suppresses useful
information.

\paragraph{Contract design.}
For instruments that do embed above-par early redemption features,
Theorem~\ref{thm:main} provides a conditional fragility result: the
withdrawal threshold rises and failure probability weakly increases
in $\varepsilon$ for survival maps satisfying (R1)--(R2). Regulatory
requirements on contract terms---floors on early redemption
penalties, restrictions on above-par early redemption in wholesale
funding---have a direct theoretical basis in the upward-slope
result. The 2010 money market fund reforms and the 2016 floating
net asset value rules for institutional funds are consistent with
this prediction.

\subsection{Empirical Predictions and Data Requirements}
\label{sec:predictions}

The model yields four design-conditional predictions that distinguish
its information channel from a fundamentals-only benchmark.
\begin{enumerate}
  \item Conditional on reserves, asset quality, deposit size, and
    insurance coverage, cascade risk varies with the depositor
    observation network.
  \item An early withdrawal by a depositor with greater pre-run
    observation centrality affects more subsequent decisions.
  \item In the present NWS specification, larger neighborhoods reduce
    local sampling variance and false inference. They do not guarantee
    unbiased beliefs when network position and withdrawal behavior are
    correlated.
  \item Delaying or batching the observation of withdrawals changes
    cascade probability even when contract payoffs and fundamentals
    are held fixed.
\end{enumerate}

Testing these predictions requires timestamped depositor withdrawals,
balances and insurance status, bank liquidity and asset-quality
controls, and pre-event links constructed from ownership,
transactions, communication, professional affiliation, or common
information channels. The economically relevant observation network
must be measured before the run rather than inferred from withdrawals
that the model seeks to explain.

\citet{cookson2023social} show that social-media activity among
SVB-connected venture-capital firms and startups temporally predicts
deposit outflows. This is consistent with information propagation,
but it neither identifies a structural network effect nor estimates
the parameters above.

The NWS graph is a tractable two-parameter baseline. Its degree $k$
governs local sample size, while shortcut density changes clustering
and path length; these are distinct objects. Real depositor networks
may also exhibit assortativity, centrality concentration, and degree
heterogeneity not captured by NWS. The present evidence therefore
supports claims about local sample size and shortcut density in the
specified designs, not topology invariance across arbitrary graph
families.

%%============================================================
\section{Conclusion}
\label{sec:conclusion}
%%============================================================

One component of the canonical Diamond--Dybvig mechanism is an
above-par premium channel. Conditional on the imposed survival map,
the withdrawal threshold rises and failure probability weakly
increases with the early-withdrawal premium $\varepsilon$; at the premium levels
in ordinary deposit contracts ($\varepsilon \lesssim 0.02$),
Corollary~\ref{cor:approx} shows how the withdrawal threshold departs
locally from its zero-premium limit, while the paired finite-agent
bridge simulation verifies that the limiting recovery comparison is
usually robust to alternative classifications of partial payment.
Ordinary demand deposits have no above-par premium, so that DD channel
is not carried unchanged into their run incentives.
\citet{diamond1983bank} proved that bank runs can occur as a pure
coordination failure---an existence result that stands. This paper
identifies how the canonical contract changes the marginal withdrawal
condition, while preserving the relevance of sequential service,
recovery priority, and strategic complementarity at par.

The core formal results are two. Conditional on the imposed survival
map, Theorem~\ref{thm:main} establishes that the withdrawal threshold
rises and failure probability weakly increases with the insurance
premium $\varepsilon=c_1-1$ embedded in the DD optimal contract,
operating through three reinforcing
channels: the option becomes more valuable, the continuation value
deteriorates, and withdrawal events accelerate pool depletion.
Theorem~\ref{thm:bridge} shows that as $\varepsilon \to 0$---the
realistic limit---rational CRRA agents reduce to probability
threshold comparators. In the finite-agent bridge simulation, the two
extreme partial-payment bounds agree in 153 of
180 scenarios. The expanded DD experiment independently confirms the
aggregate upward premium comparative static in 41{,}600
fixed-contract cells, while 7{,}840 allocation cells separately
evaluate participation and utility-shift sensitivity using fresh
holdout draws.

The deeper contribution is methodological. DD's static two-equilibrium framework made it natural to ask whether
runs could occur, and left open what determines the path between
equilibria. Stochastic process techniques make that path explicit.
When marginal agent behavior is modeled explicitly, the option
premium emerges as one structural driver of equilibrium selection
inside the canonical contract. Removing it changes the cutoff but
does not remove the underlying sequential-service problem.

The network-contagion model generates coordination failure through
information contagion on depositor networks. Agents observe local
withdrawal rates, treat them as proxies for global conditions---an
internally coherent response under the maintained belief model---and
withdraw when their
estimated loss probability crosses the threshold $p^*$ microfounded
by Theorem~\ref{thm:bridge}. High-withdrawal terminal states emerge
from cascade dynamics that depend on neighborhood degree, shortcut
density, and activation order. The matched simulations establish that local information
under sequential service is sufficient to produce endogenous
withdrawals and failures under a comparative recovery rule, the
bridge threshold, and explicit expected utility. The latter two rules
agree on 99.5\% of validation
failure outcomes and 96.1\% of
the prespecified sweep, although they diverge in deliberately selected
boundary cases. The network mechanism is therefore robust to the
precise microfoundation over most prespecified settings without
making the stronger claim that the decision rule never matters.
Unlike the canonical premium parameter, which is absent from
par-value demand deposits, reserve ratios and insurance coverage are
measurable; depositor observation links may be recoverable once their
economic definition is specified.

Deposit insurance operates on recovery risk, while information-timing
policies operate on what later agents observe. These are
model-implied intervention margins rather than welfare rankings.
The SVB evidence is compatible with rapid information propagation
\citep{cookson2023social}, but does not identify the model against
other run mechanisms.

%%============================================================
%% Appendices
%%============================================================

\appendix

\section{Expanded Diamond--Dybvig Sweep}
\label{app:ddsweep}

The expanded DD exercise contains three prespecified panels. The
fixed-contract sensitivity panel has 41{,}600 cells: eight withdrawal
premia, four productivity values, 100 prior withdrawal probabilities,
and thirteen nonredundant $(\rho,a)$ specifications. The main
allocation panel has 6{,}400 cells and fixes $a=1$ while searching
101 common deposit allocations. The targeted shift-allocation panel
adds 1{,}440 cells at $a\in\{0.1,10,100\}$ for nonlinear utility.
Risk-neutral utility is evaluated only at $a=1$ because the shift is
an additive constant when $\rho=0$ and cannot alter choices. Every
allocation candidate uses 1{,}000 search realizations, and every
selected or fixed allocation uses 10{,}000 independent holdout
realizations.

\begin{table}[ht]
\centering
\small
\caption{Fixed-contract DD outcomes by withdrawal premium}
\label{tab:ddexpanded}

\begin{tabular}{rrr}
\toprule
Withdrawal premium & Mean failure rate & Mean endogenous withdrawals\\
\midrule
0.000 & 0.000 & 0.09\\
0.005 & 0.058 & 0.11\\
0.010 & 0.065 & 0.13\\
0.020 & 0.092 & 0.26\\
0.050 & 0.169 & 0.73\\
\addlinespace
0.100 & 0.271 & 1.79\\
0.250 & 0.482 & 4.85\\
0.500 & 0.677 & 7.91\\
\bottomrule
\end{tabular}


\smallskip

\small\textit{Note}: Baseline utility shift $a=1$. Each row averages
over four risk-aversion values, four productivities, and 100 prior
withdrawal probabilities. Deposits are fixed at 1{,}000.
\end{table}

Table~\ref{tab:ddexpanded} shows the aggregate premium gradient.
Zero premium produces endogenous timing responses in some
high-withdrawal cells but no contractual shortfall: the full pool can
pay all claims at par. Positive premia create a shortfall boundary,
which moves inward as the premium rises. The full matched comparison
contains nine local reversals greater than one percentage point out
of 5{,}200 premium paths; six exceed two points. These exceptions are
reported rather than smoothed away and do not overturn the aggregate
boundary movement.

\begin{figure}[ht]
\centering
\begin{knitrout}
\definecolor{shadecolor}{rgb}{0.969, 0.969, 0.969}\color{fgcolor}
\includegraphics[width=0.88\textwidth]{figures/knitr-dd-shift-figure-1} 
\end{knitrout}
\caption{Common-allocation sensitivity to the shifted-CRRA
parameter. Participation classification is identical across all four
shifts in every matched cell. Deposit differences are concentrated
at higher risk aversion and remain small relative to the 1{,}000-unit
endowment.}
\label{fig:ddshift}
\end{figure}

\section{Proof of Proposition~\ref{prop:upslope}}
\label{app:proofA}

Proposition~\ref{prop:upslope} is a direct consequence of
Theorem~\ref{thm:main}; we prove the theorem. We work with shifted
CRRA utility throughout:
\[
  u(c) = \frac{(1+c)^{1-\rho}}{1-\rho}, \quad \rho > 0,\;\rho \neq 1;
  \qquad u(c) = \log(1+c), \quad \rho = 1.
\]
Under shifted CRRA the utility of zero consumption is $u(0) =
1/(1-\rho)$ for $\rho \neq 1$ and $u(0)=0$ for $\rho=1$.
An agent who stays earns $c_2$ if the bank is solvent and earns
nothing otherwise. Withdrawing iff $u(c_1) \geq P\,u(c_2) + (1-P)\,u(0)$
rearranges to $P \leq \Phi_{\rm shift}$, where the withdrawal
threshold is (see Appendix~\ref{app:proofB}, Step~1):
\begin{equation}
  \Phi_{\rm shift}(c_1,\, c_2,\, \rho) \;=\;
  \frac{(1+c_1)^{1-\rho} - 1}{(1+c_2)^{1-\rho} - 1}
  \quad (\rho \neq 1);
  \qquad
  \Phi_{\rm shift}(c_1,\, c_2,\, 1) \;=\; \frac{\log(1+c_1)}{\log(1+c_2)}.
  \label{eq:phishiftA}
\end{equation}
Because $u(c_2)-u(0) > 0$ for all $\rho > 0$ and $c_2 > 0$, the
division step preserves the inequality direction for all $\rho > 0$.
An agent withdraws if and only if $P(\text{survival}\mid D) \leq
\Phi_{\rm shift}$.

\medskip

\noindent\textit{Resource constraint.}\; Differentiating the
binding resource constraint~\eqref{eq:resource} at $n = \pi$
(the fundamental equilibrium fraction) gives:
\begin{equation}
  c_2(c_1) \;=\; \frac{R(1-\pi c_1)}{1-\pi},
  \qquad
  c_2'(c_1) \;=\; -\frac{\pi R}{1-\pi} \;<\; 0.
  \label{eq:c2deriv}
\end{equation}
This is \textit{Channel~2}: the continuation value is strictly
decreasing in $c_1$.

\begin{lemma}
\label{lem:phi}
For all $\rho > 0$, $\Phi_{\rm shift}\bigl(c_1,\, c_2(c_1),\, \rho\bigr)$
is strictly increasing in $c_1$ along the resource
constraint~\eqref{eq:c2deriv}.
\end{lemma}

\begin{proof}
\textit{Case $\rho \neq 1$.}\; Let $\alpha = 1-\rho \neq 0$ and write
$\Phi_{\rm shift} = f/g$ where $f \equiv (1+c_1)^\alpha - 1$ and
$g \equiv (1+c_2)^\alpha - 1$. For any $x > 0$, $(1+x)^\alpha > 1$
iff $\alpha > 0$, so $f$ and $g$ share the sign of $\alpha$ and
$g \neq 0$. By the quotient rule:
\[
  \frac{d\Phi_{\rm shift}}{dc_1}
  \;=\; \frac{f'g - fg'}{g^2}.
\]
Substituting $f' = \alpha(1+c_1)^{\alpha-1}$ and
$g' = \alpha(1+c_2)^{\alpha-1}c_2'(c_1)$ with
$c_2'(c_1) = -|c_2'| < 0$:
\[
  f'g - fg' \;=\;
  \alpha\!\left[\,(1+c_1)^{\alpha-1}\,g
               \;+\; (1+c_2)^{\alpha-1}|c_2'|\,f\,\right].
\]
Each of $(1+c_1)^{\alpha-1}$ and $(1+c_2)^{\alpha-1}$ is strictly
positive. Both $f$ and $g$ share the sign of $\alpha$, so each
summand in the bracket shares the sign of $\alpha$, and hence the
bracket itself shares the sign of $\alpha$. Multiplying by $\alpha$
gives a product proportional to $\alpha^2 > 0$. Since $g^2 > 0$,
we conclude $d\Phi_{\rm shift}/dc_1 > 0$.

\textit{Case $\rho = 1$.}\; Write $\Phi_{\rm shift} = f/g$ with
$f = \log(1+c_1)$ and $g = \log(1+c_2(c_1))$. Then:
\[
  \frac{d\Phi_{\rm shift}}{dc_1}
  \;=\; \frac{f'g - fg'}{g^2}
  \;=\; \frac{1}{g^2}
    \left[\frac{g}{1+c_1}
      - \log(1+c_1)\cdot\frac{c_2'(c_1)}{1+c_2}\right].
\]
Both terms are positive ($g > 0$ since $c_2 > 0$; the second term
is positive because $c_2'(c_1) < 0$). Hence $d\Phi_{\rm shift}/dc_1> 0$.
\end{proof}

\begin{lemma}
\label{lem:dstar}
The rational expectations equilibrium threshold $D^*$ is strictly
increasing in $c_1$.
\end{lemma}

\begin{proof}
The equilibrium condition is that an agent activated when
$D = D^*$ is exactly indifferent:
$P(\text{survival to }T\mid D(0)=D^*)=\Phi_{\rm shift}$.
Define $G(D^*,\, c_1) \equiv
P(\text{survival to }T \mid D(0) = D^*) -
\Phi_{\rm shift}(c_1,\, c_2(c_1),\, \rho)$, so equilibrium
requires $G = 0$.

Write the survival probability more fully as
$S(D;\,D^*,c_1)$. Regularity condition~(R1) applies to the
equilibrium composite $S(D^*;\,D^*,c_1)$ and gives
$\partial G/\partial D^*>0$: a larger equilibrium pool raises the
probability of surviving through maturity. Survival also
depends directly on $c_1$ because $c_1$ is the size of each
withdrawal payment. Holding $D$ and the withdrawal region fixed, the
coupling used in Lemma~\ref{lem:jump} gives
$\partial S/\partial c_1\leq0$. By Lemma~\ref{lem:phi},
$d\Phi_{\rm shift}/dc_1>0$, and therefore
\[
  \frac{\partial G}{\partial c_1}
  =
  \frac{\partial S}{\partial c_1}
  -
  \frac{d\Phi_{\rm shift}}{dc_1}
  <0.
\]
The implicit function theorem gives:
\[
  \frac{dD^*}{dc_1}
  \;=\; -\frac{\partial G/\partial c_1}{\partial G/\partial D^*}
  \;>\; 0.
\]
Since agents withdraw when $D \leq D^*$ (the pool level at which
$P(\text{survival}\mid D) \leq \Phi_{\rm shift}$), a higher $D^*$
strictly enlarges the withdrawal region. \qed
\end{proof}

\begin{lemma}
\label{lem:jump}
For any fixed $D^*$, the failure probability $P(\tau < T)$ is
weakly increasing in $c_1$. It is strictly increasing if a
positive-probability set of coupled paths survives under $c_1$ but
fails under $c_1'>c_1$ (Channel~3: jump-size effect).
\end{lemma}

\begin{proof}
Consider two economies with identical $D^*$ and $D(0)$ but
withdrawal payments $c_1 < c_1'$. Couple the two processes on the
same probability space, using the same exogenous withdrawal
realizations and the same activation-order draw. Each withdrawal
event reduces the deposit pool by $c_1$ in the first economy and
$c_1' > c_1$ in the second. Pathwise, $D'(t) \leq D(t)$ for every
$t$: any trajectory that reaches $D(t) = 0$ under $c_1$ also
reaches $D'(t) = 0$ under $c_1'$. Hence $P(\tau' < T) \geq
P(\tau < T)$. Strict inequality follows only when trajectories that
survive under $c_1$ but reach $D'(t)=0$ under $c_1'$ have positive
probability. Without such boundary-crossing paths, the weak
inequality may hold with equality even away from unit failure
probability.
\end{proof}

\begin{proof}[Proof of Theorem~\ref{thm:main}]
Lemma~\ref{lem:dstar} shows $D^*$ is strictly increasing in $c_1$.
A higher $D^*$ enlarges the set of deposit-pool states from which
agents withdraw, weakly increasing $P(\tau < T)$. Lemma~\ref{lem:jump}
shows $P(\tau < T)$ is further weakly increased by the larger
per-withdrawal jump size, for any fixed $D^*$, with strict
inequality under the positive-probability crossing condition. Both effects operate
simultaneously and in the same direction. On the interior region
$\{c_1 : 0 < P(\tau < T) < 1\}$ at least one effect is strict,
establishing strict monotonicity. At $P(\tau < T) = 1$ the result
holds with weak inequality only. Proposition~\ref{prop:upslope}
restates the interior result.
\end{proof}

\begin{remark}[Extension to partial connectivity]
\label{rem:partialconn}
The proof constructs a single equilibrium threshold $D^*$,
which corresponds to the case where all agents observe a common
signal---the full deposit pool, as under complete network
connectivity. Under partial connectivity each agent $i$ observes
a neighborhood signal $D_i$ drawn from $m_i < K$ peers, so agents
hold heterogeneous beliefs simultaneously and no single $D^*$
characterizes all withdrawal decisions at once.

The result extends agent by agent under the model's belief
specification. The posterior $p_S(d)$ is computed from the
exogenous withdrawal probability $q = 1 - p_0$---the prior
probability that any observed withdrawal is an exogenous
shock---not from the endogenous cascade strategies of other
agents. Consequently $p_S(d)$ does not depend on other agents'
thresholds $D^*_{-i}$, the indifference condition for agent $i$
is determined by $(c_1,\, c_2,\, \rho,\, k_i)$ alone, and the
implicit-function argument of Lemma~\ref{lem:dstar} applies to
each $D^*_i$ separately, giving $dD^*_i/dc_1 > 0$ for every $i$.

If beliefs are endogenized to the cascade, a lattice argument is
available only under an additional exact-monotonicity restriction.
Specifically, suppose the strategy map $T^0$ evaluates probabilities
without sampling error, fixes insurance and other bank-state
components, and has the property that raising any other agent's
withdrawal threshold weakly lowers $p_S(d)$ at every $d$. Then $T^0$
is isotone on the finite strategy lattice, so Tarski's theorem yields
minimal and maximal fixed points and standard monotone comparative
statics apply.

These restrictions define an analytical threshold benchmark. They
are not asserted for every realized network simulation, whose finite
Monte Carlo estimates, reserve depletion, partial payments, and
changing insurance states can break realized isotonicity. The
network evidence therefore concerns terminal states of the
implemented irreversible cascade process; it is not presented as a
computational verification of the Tarski benchmark.
\end{remark}

\begin{remark}[Partial payout and the large-$K$ limit]
\label{rem:partialpayout}
The withdrawal condition~\eqref{eq:wdcond} treats the failure payoff
as $u(0)$, the utility of zero consumption. In the finite-$K$
simulation the bank pays $\min(c_1,\,\text{vault})$, so exactly
one agent---the last before vault exhaustion---may receive a
partial amount $v = \text{vault} \bmod\, c_1 \in (0, c_1)$ rather
than zero. This does not affect Proposition~\ref{prop:upslope} in
the large-$K$ limit. For any $K$-agent economy the probability
that a given agent is the partial-payout agent is exactly $1/K$,
so the correction to the withdrawal threshold $\Phi_{\rm shift}$
is at most $\Phi_{\rm shift}/K \leq 1/K$ in absolute value---an
$O(1/K)$ perturbation that vanishes as $K \to \infty$. To see
this, let $\Phi_v = [u(v)-u(0)]/[u(c_2)-u(0)]$ be the threshold
an agent would use if the failure payoff were $v$ with certainty.
Then $\Phi_{\rm partial} = (1-1/K)\Phi_{\rm shift} + (1/K)\Phi_v$
and $|\Phi_{\rm partial}-\Phi_{\rm shift}| = (\Phi_{\rm
shift}-\Phi_v)/K \leq \Phi_{\rm shift}/K < 1/K$, using $\Phi_v
\geq 0$ and regularity condition~(R2). At $K = 1{,}000$ (the
simulation size) this correction is of order $10^{-3}$ and does
not affect the qualitative monotonicity. The residual $v =
\text{vault} \bmod\, c_1$ is non-monotone in $c_1$ in general,
so no clean first-order stochastic dominance argument closes the
finite-$K$ case; Proposition~\ref{prop:upslope} is therefore
stated as a large-$K$ result, confirmed numerically for the finite
model in Section~\ref{sec:ddsimulation}.
\end{remark}

\section{Proof of Theorem~\ref{thm:bridge} (Bridge Theorem)}
\label{app:proofB}

We prove Theorem~\ref{thm:bridge}. The proof has three steps:
writing down the withdrawal condition under shifted CRRA with the
correct outside option, taking the limit $\varepsilon \to 0$, and
establishing continuity at $\rho = 1$ by L'H\^opital's rule.

\medskip

\noindent\textit{Step 1: full withdrawal condition.}

With shifted CRRA, $u(c) = (1+c)^{1-\rho}/(1-\rho)$ for $\rho
\neq 1$ and $u(c) = \log(1+c)$ for $\rho = 1$, the utility of
consuming zero is $u(0) = 1/(1-\rho)$ for $\rho \neq 1$ and
$u(0) = 0$ for $\rho = 1$. A depositor who stays earns $c_2$ if
the bank remains solvent and earns nothing otherwise. Expected
utility of staying is:
\[
  \mathbb{E}[u \mid \text{stay}]
  \;=\; P\,u(c_2) \;+\; (1-P)\,u(0).
\]
The agent withdraws iff $u(c_1) \geq \mathbb{E}[u \mid \text{stay}]$,
i.e., $u(c_1) - u(0) \geq P\,[u(c_2) - u(0)]$, so the withdrawal
threshold is:
\begin{equation}
  P \;\leq\; \Phi_{\rm shift}(\varepsilon)
  \;\equiv\;
  \frac{u(c_1) - u(0)}{u(c_2) - u(0)}
  \;=\;
  \frac{(1+c_1)^{1-\rho} - 1}{(1+c_2)^{1-\rho} - 1}
  \quad (\rho \neq 1).
  \label{eq:phishift}
\end{equation}
For $\rho = 1$: $\Phi_{\rm shift}(\varepsilon) =
\log(1+c_1)/\log(1+c_2)$.

\medskip

\noindent\textit{Step 2: limit $\varepsilon \to 0$.}

Set $c_1(\varepsilon) = 1 + \varepsilon$. The resource
constraint~\eqref{eq:resource} at fundamental withdrawal fraction
$n = \pi$ gives $c_2(\varepsilon) = R(1 - \pi c_1)/(1-\pi)$,
which by equation~\eqref{eq:c2eps} satisfies $c_2(0) = R$ and
$c_2'(0) = -\pi R/(1-\pi) < 0$. As $\varepsilon \to 0$, $c_1 \to
1$ and $c_2 \to R$. By continuity of $\Phi_{\rm shift}$ in
$(c_1, c_2)$:
\[
  \lim_{\varepsilon \to 0}\,\Phi_{\rm shift}(\varepsilon)
  \;=\; \frac{2^{1-\rho} - 1}{(1+R)^{1-\rho} - 1}
  \;\equiv\; p^*(\rho,\, R)
  \quad (\rho \neq 1).
\]
For $\rho = 1$:
$\lim_{\varepsilon \to 0}\,\Phi_{\rm shift}(\varepsilon) =
\log(2)/\log(1+R) \equiv p^*(1,\,R)$.
The withdrawal condition therefore converges to $P \leq p^*(\rho,R)$
as stated in Theorem~\ref{thm:bridge}.

\medskip

\noindent\textit{Step 3: continuity at $\rho = 1$.}

Write $x = 1 - \rho \to 0$:
\[
  p^*(\rho,\,R) \;=\; \frac{2^x - 1}{(1+R)^x - 1}.
\]
Both numerator and denominator tend to zero as $x \to 0$.
L'H\^opital's rule gives:
\[
  \lim_{x \to 0}\,\frac{2^x - 1}{(1+R)^x - 1}
  \;=\; \lim_{x \to 0}\,\frac{2^x \log 2}{(1+R)^x \log(1+R)}
  \;=\; \frac{\log 2}{\log(1+R)}
  \;=\; p^*(1,\,R).
\]
Hence $p^*$ is continuous at $\rho = 1$.

\medskip

\noindent\textit{Remark.} Since $R > 1$, we have $1 + R > 2$, so
$(1+R)^{1-\rho} > 2^{1-\rho}$ for $\rho < 1$, giving $p^* < 1$.
Clearly $p^* > 0$. Hence the threshold is a proper probability
for all $\rho \in (0,1]$ and $R > 1$. \qed

\begin{proof}[Proof of Corollary~\ref{cor:approx}]
Since $\Phi_{\rm shift}(\varepsilon) = [(2+\varepsilon)^\alpha - 1]\,/\,
[(1+c_2(\varepsilon))^\alpha - 1]$ with $c_2(\varepsilon)$
smooth and $c_2(0) = R$, the function $\Phi_{\rm shift}$ is $C^1$
at $\varepsilon = 0$ (the denominator is non-zero there). The
Taylor expansion~\eqref{eq:approx} therefore holds with slope
$M = d\Phi_{\rm shift}/d\varepsilon\big|_{\varepsilon=0}$.

To compute $M$, differentiate using the quotient rule with
$f(\varepsilon) = (2+\varepsilon)^\alpha - 1$ and
$g(\varepsilon) = (1+c_2(\varepsilon))^\alpha - 1$:
\[
  M = \left.\frac{f'g - fg'}{g^2}\right|_{\varepsilon=0}.
\]
At $\varepsilon = 0$: $f(0)=B$, $g(0)=A$,
$f'(0) = \alpha\cdot 2^{\alpha-1}$, and
$g'(0) = \alpha(1+R)^{\alpha-1}c_2'(0)=-\alpha\mu(1+R)^{\alpha-1}$
where $\mu = \pi R/(1-\pi)$. Substituting:
\[
  M = \frac{\alpha\cdot 2^{\alpha-1}\cdot A
            \;+\; \alpha\mu(1+R)^{\alpha-1}\cdot B}{A^2}
    = \frac{\alpha\bigl[2^{\alpha-1}A + \mu(1+R)^{\alpha-1}B\bigr]}{A^2},
\]
which is equation~\eqref{eq:Mformula}. The sign $M > 0$
follows from the $\alpha^2$ argument: $A$ and $B$ share the
sign of $\alpha$, both power factors are positive, so the
bracket shares the sign of $\alpha$, and $\alpha\cdot
\mathrm{sign}(\alpha) = \alpha^2 > 0$. The case $\rho = 1$
follows by differentiating $\Phi_{\rm shift}(\varepsilon) =
\log(2+\varepsilon)/\log(1+R-\mu\varepsilon)$ directly at
$\varepsilon = 0$, yielding equation~\eqref{eq:Mformula_log}.
\end{proof}

\section{Agent-Based Model: Implementation Details}
\label{app:abm}

\subsection{Frozen 2026 Simulation Design}

The production simulations are implemented in Julia 1.11.3. Every
stage records stable seeds, Julia version, execution representation,
and immutable design markers. Runs are written to resumable chunks,
merged only after job-identifier checks, and frozen before subsequent
selection or analysis. The final integrity audit contains 96{,}000
unique confirmatory jobs and 39{,}600 unique monotonicity jobs, with
no missing analysis values and no above-par payments.

\paragraph{Diamond--Dybvig panels and bridge.}
The DD experiment uses homogeneous primitives and random sequential
activation. Its 41{,}600-cell fixed-contract panel holds deposits at
1{,}000 and crosses eight premia, four productivities, 100 prior
withdrawal probabilities, and thirteen nonredundant
risk-aversion/utility-shift specifications. Each cell uses 10{,}000
realizations. The 6{,}400-cell main allocation panel fixes the shift
at one and searches 101 allocations using 1{,}000 realizations per
candidate before evaluating the selected allocation on 10{,}000
independent holdout realizations. The 1{,}440-cell targeted shift
panel repeats selected allocation comparisons at shifts 0.1, 10, and
100. Common random numbers are shared across premium and shift
comparisons within matched fundamental groups.

The finite-agent bridge is a distinct model rather than a network
special case. For each of 180 scenarios, it uses 10{,}000 common
withdraw/stay trial draws. The payment record distinguishes full,
partial, and zero recovery for both actions. The two reported bounds
reverse the classification of partial payments: the
withdrawal-favoring bound treats partial withdrawal as success and
partial staying as failure, while the staying-favoring bound does the
opposite.

\paragraph{Network production design.}
The network model has 200 agents per production scenario and 500
paired decision draws per activation. Deposits may be homogeneous,
clipped LogNormal, or clipped Pareto. Topologies include complete,
Watts--Strogatz, Erdős--Rényi, and Barabási--Albert graphs. Insurance
may be absent, fixed, quantile-based, or adaptive. Exogenous shocks
vary in size and may be randomly or locally positioned.

Agents observe withdrawn neighbors but not their global queue
positions. A subjective state is rebuilt after every objective
withdrawal, and agents who have already withdrawn are excluded from
future simulated orders. The withdraw-now and stay actions use common
random draws. Bank payments are capped at par under every rule.

Three immutable decision specifications are run with matched scenario
and model seeds: comparative recovery probability, a bridge threshold
calibrated at 0.80, and explicit shifted-CRRA utility with $\rho=1$
and gross late return 1.25. The fixed 0.80 equals $1/1.25$ in the
binary-recovery risk-neutral calibration. It is a reduced-form member
of the cutoff family established by Theorem~\ref{thm:bridge}, not the
theorem's $p^*$ for every utility, payoff, or deposit calibration. The explicit
utility specification is consequently a robustness comparison, not a
second label for an algebraically identical numerical rule. The
1{,}800-run validation uses 600 common scenarios.
The confirmatory stage uses 80 adaptively selected structural cells,
400 replications, and three rules, for 96{,}000 jobs. The
outcome-independent monotonicity design uses six structures, eleven
settings, 200 replications, and three rules, for 39{,}600 jobs.
Prespecified paths vary reserve ratios, fixed-insurance payments, and
shock sizes.

\subsection{Adaptive Boundary Stress Test}
\label{app:adaptive}

The smart parameter sweep is deliberately separated from the
prespecified evidence in the main text because its cells are selected
using earlier simulation outcomes. A broad coverage stage first maps
the heterogeneous deposit, topology, insurance, reserve, shock-size,
and shock-location grid. A scoring rule then prioritizes cells with
high cross-rule entropy, disagreement, and statistical uncertainty,
subject to diversity constraints across economic categories. The
adaptive batch evaluates 100 selected cells under all three decision
rules with 500 replications per rule, for 150{,}000 runs.

The final selector relearns the scores from that high-replication
adaptive batch and freezes 80 cells for confirmatory execution.
Each selected cell is evaluated under the comparative,
bridge-threshold, and explicit-utility rules with 400 new
replications, producing 96{,}000 jobs. Selection and execution use
separate stable seeds. The selection file, scores, source
fingerprints, and frozen marker are retained with the run.

Because this design enriches disagreement rather than representing
the parameter grid, it is a stress test of the decision-rule
robustness claim. In the selected sample, bridge-threshold and
explicit-utility failure outcomes agree in
72.7\% of matched
cell-replications. Their failure rates are
74.9\% and
56.1\%, a difference of
18.8\%. These estimates
show that the exact decision rule matters materially in environments
chosen because the rules are likely to disagree. They do not estimate
an unconditional or population-average decision-rule effect. For
that reason, the paper's primary robustness evidence comes from the
common-scenario validation and the outcome-independent monotonicity
design reported in Section~\ref{sec:results}.

\begin{figure}[ht]
\centering
\begin{knitrout}
\definecolor{shadecolor}{rgb}{0.969, 0.969, 0.969}\color{fgcolor}
\includegraphics[width=0.86\textwidth]{figures/knitr-decision-rule-robustness-figure-1} 
\end{knitrout}
\caption{Decision-rule robustness across designs. Agreement is high
in the compact validation grid and the outcome-independent
prespecified sweep, but lower in the adaptively selected stress test.
Because the designs target different regions, the bars are
design-conditional summaries rather than estimates from a common
population.}
\label{fig:decision_rule_robustness}
\end{figure}

The large sparse execution representation shares objective deposits
and stores sparse withdrawal indices rather than cloning whole
subjective networks. Exact parity tests compare it with the small
object representation under all three economic rules. The full Julia
test suite and frozen-result integrity checks pass before the knitr
analysis reads the files.

% Legacy empirical block 2 archived in archive/legacy_empirical_blocks_20260724.tex.

%%============================================================
\bibliography{references}
%%============================================================

\end{document}
